\documentclass[11pt]{article}
\usepackage[a4paper,margin=0.65in]{geometry}
\usepackage[T1]{fontenc}
\usepackage{amsmath,amssymb,amsfonts}
\usepackage{graphicx}
\usepackage{pdfpages}
\usepackage[english,shorthands=off]{babel}
\usepackage{fancyhdr}
\usepackage[bookmarks=false,colorlinks=false]{hyperref}
\usepackage[protrusion=true,expansion=false]{microtype}
\emergencystretch=3em
\tolerance=600
\providecommand{\modd}{\mathrm{mod}}
\providecommand{\Disc}{\mathrm{Disc}}
\providecommand{\canont}{}
\providecommand{\asterism}{*}
\providecommand{\Hessian}{H}
\providecommand{\tome}{}
\providecommand{\page}{p.}
\pagestyle{fancy}\fancyhf{}\fancyhead[L]{Pages 251--300}\fancyhead[R]{Cayley --- Collected Papers, Vol. II}\renewcommand{\headrulewidth}{0.4pt}
\setlength{\headheight}{15pt}

\begin{document}

%==============================================================
% Paper 140 (PDF p. 251 onward)
%==============================================================
\begin{center}
{\Large\bfseries 140.}\\[1ex]
{\large\bfseries RESEARCHES ON THE PARTITION OF NUMBERS.}
\end{center}

\noindent\textit{[From the Philosophical Transactions of the Royal Society of London, vol.~CXLV for the year 1855, pp.~127--140. Received April 14,---Read May 24, 1855.]}

\medskip

\noindent I \textsc{propose} to discuss the following problem: ``To find in how many ways a number $q$ can be made up of the elements $a$, $b$, $c$, \ldots\ each element being repeatable an indefinite number of times.'' The required number of partitions is represented by the notation
\[
P(a,\, b,\, c,\,\ldots)\,q,
\]
and we have, as is well known,
\[
P(a,\, b,\, c,\,\ldots)\,q = \text{coefficient } x^q \text{ in } \frac{1}{(1-x^a)(1-x^b)(1-x^c)\cdots},
\]
where the expansion is to be effected in ascending powers of $x$.

It may be as well to remark that each element is to be considered as a separate and distinct element, notwithstanding any equalities which may exist between the numbers $a$, $b$, $c$, \ldots; thus, although $a=b$, yet $a+a+a+\&\text{c.}$ and $a+a+b+\&\text{c.}$ are to be considered as two different partitions of the number $q$, and so in all similar cases.

The solution of the problem is thus seen to depend upon the theory, to which I now proceed, of the expansion of algebraical fractions.

Consider an algebraical fraction $\dfrac{\phi x}{f x}$, where the denominator is the product of any number of factors (the same or different) of the form $1-x^m$. Suppose in general that $[1-x^m]$ denotes the irreducible factor of $1-x^m$, i.e.\ the factor which, equated to zero, gives the prime roots of the equation $1-x^m=0$. We have
\[
1-x^m = \prod [1-x^{m'}],
\]
where $m'$ denotes any divisor whatever of $m$ (unity and the number $m$ itself not excluded). Hence, if $a$ represent a divisor of one or more of the indices $m$, and $k$ be the number of the indices of which $a$ is a divisor, we have
\[
f x = \prod [1-x^a]^k.
\]

Now considering apart from the others one of the multiple factors $[1-x^a]^k$, we may write $f x = [1-x^a]^k f_1 x$.

Suppose that the fraction $\dfrac{\phi x}{f x}$ is decomposed into simpler fractions, in the form
\[
\frac{\phi x}{f x} = I(x) + (x \partial_x)^{k-1}\frac{\theta x}{[1-x^a]} + (x \partial_x)^{k-2}\frac{\theta_1 x}{[1-x^a]} + \cdots + \frac{\theta_{k-1} x}{[1-x^a]} + \&\text{c.},
\]
where $I(x)$ denotes the integral part, and the $\&\text{c.}$ refers to the fractional terms depending upon the other multiple factors such as $[1-x^b]^{k_1}$. The functions $\theta x$ are to be considered as functions with indeterminate coefficients, the degree of each such function being inferior by unity to that of the corresponding denominator; and it is proper to remark that the number of the indeterminate coefficients in all the functions $\theta x$ together is equal to the degree of the denominator $f x$.

The term $(x \partial_x)^{k-1}\dfrac{\theta x}{[1-x^a]}$ may be reduced to the form
\[
\frac{a x}{[1-x^a]^k} + \frac{g_1 x}{[1-x^a]^{k-1}} + \&\text{c.},
\]
the functions $g x$ being of the same degree as $\theta x$, and the coefficients of these functions being linearly connected with those of the function $\theta x$. The first of the foregoing terms is the only term on the right-hand side which contains the denominator $[1-x^a]^k$; hence, multiplying by this denominator and then writing $[1-x^a]=0$, we find
\[
\frac{\phi x}{f_1 x} = g x,
\]
which is true when $x$ is any root whatever of the equation $[1-x^a]=0$. Now by means of the equation $[1-x^a]=0$, $\dfrac{\phi x}{f_1 x}$ may be expressed in the form of a rational and integral function $G x$, the degree of which is less by unity than that of $[1-x^a]$. We have therefore $G x = g x$, an equation which is satisfied by each root of $[1-x^a]=0$, and which is therefore an identical equation; $g x$ is thus determined, and the coefficients of $\theta x$ being linear functions of those of $g x$, the function $\theta x$ may be considered as determined. And this being so, the function
\[
\frac{\phi x}{f x} - (x \partial_x)^{k-1}\frac{\theta x}{[1-x^a]}
\]
will be a fraction the denominator of which does not contain any power of $[1-x^a]$ higher than $[1-x^a]^{k-1}$; and therefore $\theta_1 x$ can be found in the same way as $\theta x$, and similarly $\theta_2 x$, and so on. And the fractional parts being determined, the integral part may be found by subtracting from $\dfrac{\phi x}{f x}$ the sum of the fractional parts, so that the fraction $\dfrac{\phi x}{f x}$ can by a direct process be decomposed in the above-mentioned form.

Particular terms in the decomposition of certain fractions may be obtained with great facility. Thus $m$ being a prime number, assume
\[
\frac{1}{(1-x)(1-x^2)\cdots(1-x^{m-1})(1-x^m)} = \&\text{c.} + \frac{\theta x}{[1-x^m]};
\]
then observing that $(1-x^m) = (1-x)[1-x^m]$, we have for $[1-x^m]=0$,
\[
\theta x = \frac{1}{(1-x)(1-x^2)\cdots(1-x^{m-1})}.
\]
Now $u$ being any quantity whatever and $x$ being a root of $[1-x^m]=0$, we have identically
\[
[1-u^m] = (u-x)(u-x^2)\cdots(u-x^{m-1});
\]
and therefore putting $u=1$, we have $m = (1-x)(1-x^2)\cdots(1-x^{m-1})$, and therefore
\[
\theta x = \frac{1}{m},
\]
whence
\[
\frac{1}{(1-x)(1-x^2)\cdots(1-x^m)} = \&\text{c.} + \frac{1}{m[1-x^m]}.
\]

Again, $m$ being as before a prime number, assume
\[
\frac{1}{(1-x)(1-x^2)\cdots(1-x^{m-1})(1-x^m)^2} = \&\text{c.} + \frac{\theta x}{[1-x^m]^2};
\]
we have in this case for $[1-x^m]=0$,
\[
\theta x = \frac{1}{(1-x)^2(1-x^2)\cdots(1-x^{m-1})},
\]
which is immediately reduced to $\theta x = \dfrac{1}{m^2}\,\dfrac{1}{1-x}$. Now
\[
\frac{1-u^m}{u-x} = (1+x+\cdots+x^{m-1}) + (1+x+\cdots+x^{m-2})x \cdots (1+u)x^{m-1} + x^m,
\]
or putting $u=1$,
\[
\frac{m}{1-x} = (m-1) + (m-2)x + \cdots + x^{m-1};
\]
and substituting this in the value of $\theta x$, we find
\[
\frac{1}{(1-x)(1-x^2)\cdots(1-x^m)^2} = \&\text{c.} + \frac{1}{m^3}\,\frac{(m-1)+(m-2)x+\cdots+x^{m-1}}{[1-x^m]^2}.
\]

The preceding decomposition of the fraction $\dfrac{\phi x}{f x}$ gives very readily the expansion of the fraction in ascending powers of $x$. For, consider a fraction such as
\[
\frac{\theta x}{[1-x^a]}
\]
where the degree of the numerator is in general less by unity than that of the denominator; we have
\[
1-x^a = [1-x^a]\,\prod[1-x^{a'}],
\]
where $a'$ denotes any divisor of $a$ (including unity, but not including the number $a$ itself). The fraction may therefore be written under the form
\[
\frac{\theta x\,\prod[1-x^{a'}]}{1-x^a},
\]
where the degree of the numerator is in general less by unity than that of the denominator, i.e.\ is equal to $a-1$. Suppose that $b$ is any divisor of $a$ (including unity, but not including the number $a$ itself), then $1-x^b$ is a divisor of $\prod[1-x^{a'}]$, and therefore of the numerator of the fraction. Hence representing this numerator by
\[
A_0 + A_1 x + \cdots + A_{a-1}x^{a-1},
\]
and putting $a = bc$, we have (corresponding to the case $b=1$)
\[
A_0 + A_1 + A_2 + \cdots + A_{a-1} = 0,
\]
and generally for the divisor $b$,
\[
A_0 + A_b + \cdots + A_{(c-1)b} = 0,
\]
\[
A_1 + A_{b+1} + \cdots + A_{(c-1)b+1} = 0,
\]
\[
A_{b-1} + A_{2b-1} + \cdots + A_{cb-1} = 0.
\]

Suppose now that $\alpha_q$ denotes a circulating element to the period $a$, i.e.\ write a function such as
\[
\alpha_q = 1, \quad q \equiv 0 \pmod{a},
\]
\[
\alpha_q = 0 \text{ in every other case};
\]
\[
A_0 \alpha_q + A_1 \alpha_{q-1} + \cdots + A_{a-1}\alpha_{q-a+1}
\]
will be a circulating function, or circulator to the period $a$, and may be represented by the notation
\[
(A_0,\, A_1,\, \ldots,\, A_{a-1})\,\text{circlor }a_q.
\]

In the case however where the coefficients $A$ satisfy, for each divisor $b$ of the number $a$, the above-mentioned equations, the circulating function is what I call a prime circulator, and I represent it by the notation
\[
(A_0,\, A_1,\, \ldots,\, A_{a-1})\,\text{per }a_q.
\]
By means of this notation we have at once
\[
\text{coefficient } x^q \text{ in } \frac{\theta x}{[1-x^a]} = (A_0,\, A_1,\, \ldots,\, A_{a-1})\,\text{per }a_q,
\]
and thence also
\[
\text{coefficient } x^q \text{ in } (x \partial_x)^r \frac{\theta x}{[1-x^a]} = q^r\,(A_0,\, A_1,\, \ldots,\, A_{a-1})\,\text{per }a_q.
\]

Hence assuming that in the fraction $\dfrac{\phi x}{f x}$ the degree of the numerator is less than that of the denominator (so that there is not any integral part), we have
\[
\text{coefficient } x^q \text{ in } \frac{\phi x}{f x} = \Sigma\,q^r\,(A_0,\, A_1,\, \ldots,\, A_{a-1})\,\text{per }a_q;
\]
or, if we wish to put in evidence the non-circulating part arising from the divisor $a=1$,
\[
\text{coefficient } x^q \text{ in } \frac{\phi x}{f x} = A q^{k-1} + B q^{k-2} + \cdots + L q + M + \Sigma\,q^r\,(A_0,\, A_1,\, \ldots,\, A_{a-1})\,\text{per }a_q;
\]
where $k$ denotes the number of the factors $1-x^m$ in the denominator $f x$, $a$ is any divisor (unity excluded) of one or more of the indices $m$; and for each value of $a$ the index $r$ extends from $r=0$ to $r=k-1$, where $k$ denotes the number of indices of which $a$ is a divisor. The particular results previously obtained show, that $m$ being a prime number,
\[
\text{coefficient } x^q \text{ in } \frac{1}{(1-x)(1-x^2)\cdots(1-x^m)} = \&\text{c.} + \frac{1}{m}(1, 0, 0, \ldots, 0)\,\text{per }m_q,
\]
and
\[
\text{coefficient } x^q \text{ in } \frac{1}{(1-x)(1-x^2)\cdots(1-x^m)^2} = \&\text{c.} + \frac{1}{m^3}(m-1,\, m-2,\, \ldots,\, 1)\,\text{per }m_q.
\]

Suppose, as before, that the degree of $\phi x$ is less than that of $f x$, and let the analytical expression above obtained for the coefficient of $x^q$ in the expansion in ascending powers of $x$ of the fraction $\dfrac{\phi x}{f x}$ be represented by $F q$, it is very remarkable that if we expand $\dfrac{\phi x}{f x}$ in descending powers of $x$, then the coefficient of $x^q$ in this new expansion ($q$ is here of course negative, since the expansion contains only negative powers of $x$) is precisely equal to $-F q$; this is in fact at once seen to be the case with respect to each of the partial fractions into which $\dfrac{\phi x}{f x}$ has been decomposed, and it is consequently the case with respect to the fraction itself\footnote{The property is a fundamental one in the general theory of developments.}. This gives rise to a result of some importance. Suppose that $\phi x$ and $f x$ are respectively of the degrees $N$ and $D$; it is clear from the form of $f x$ that we have $f\!\left(\dfrac{1}{x}\right) = (-)^D x^{-D} f x$; and I suppose that $\phi x$ is also such that $\phi\!\left(\dfrac{1}{x}\right) = (+)^N x^{-N} \phi x$; then writing $D - N = h$, and supposing that $\dfrac{\phi x}{f x}$ is expanded in descending powers of $x$, so that the coefficient of $x^q$ in the expansion is $F q$, it is in the first place clear that the expansion will commence with the term $x^{-h}$, and we must therefore have
\[
F q = 0
\]
for all values of $q$ from $q = -1$ to $q = -(h-1)$.

Consider next the coefficient of a term $x^{-h-q}$, where $q$ is $0$ or positive; the coefficient in question, the value of which is $-F(-h-q)$, is obviously equal to the coefficient of $x^{h+q}$ in the expansion in ascending powers of $x$ of $\dfrac{\phi(1/x)}{f(1/x)}$, i.e.\ to
\[
\frac{(+)^N}{(-)^D}\,\text{coefficient } x^{h+q} \text{ in } \frac{\phi x}{f x}\cdot x^{D-N},
\]
or what is the same thing, to
\[
\frac{(+)^N}{(-)^D}\,\text{coefficient } x^q \text{ in } \frac{\phi x}{f x};
\]
and we have therefore, $q$ being zero or positive,
\[
F(-h-q) = -(+)^N(-)^D F q.
\]
In particular, when $\phi x = 1$, $F q = 0$ for all values of $q$ from $q = -1$ to $q = -(D-1)$; and $q$ being $0$ or positive,
\[
F(-D-q) = (-)^{D-1} F q.
\]

The preceding investigations show the general form of the function $P(a, b, c, \ldots) q$, viz.\ that
\[
P(a, b, c, \ldots) q = A q^{k-1} + B q^{k-2} + \cdots + L q + M + \Sigma\,q^r\,(A_0, A_1, \ldots, A_{l-1})\,\text{per } l_q,
\]
a formula in which $k$ denotes the number of the elements $a$, $b$, $c$, \ldots, and $l$ is any divisor (unity excluded) of one or more of these elements; the summation in the case of each such divisor extends from $r=0$ to $r=k-1$, where $k$ is the number of the elements $a$, $b$, $c$, \ldots\ of which $l$ is a divisor; and the investigations indicate how the values of the coefficients $A$ of the prime circulators are to be obtained. It has been moreover in effect shown, that if $D = a + b + c + \cdots$, then, writing for shortness $P(q)$ instead of $P(a, b, c, \ldots) q$, we have
\[
P(q) = 0
\]
for all values of $q$ from $q = -1$ to $q = -(D-1)$, and that $q$ being $0$ or positive,
\[
P(-D - q) = (-)^{k-1} P(q);
\]
these last theorems are however uninterpretable in the theory of partitions, and they apply only to the analytical expression for $P(q)$.

I have calculated the following particular results:
\begin{align*}
P(1, 2)q &= \tfrac{1}{4}\Bigl\{2q + 3 + (1,\,-1)\,\text{per }2_q\Bigr\}, \\[1ex]
P(1, 2, 3)q &= \tfrac{1}{72}\Bigl\{6q^2 + 36q + 47 \\
&\qquad\quad + 9(1,\,-1)\,\text{per }2_q \\
&\qquad\quad + 8(2,\,-1,\,-1)\,\text{per }3_q\Bigr\}, \\[1ex]
P(1, 2, 3, 4)q &= \tfrac{1}{288}\Bigl\{2 q^3 + 30 q^2 + 135 q + 175 \\
&\qquad\quad + (9q + 45)(1,\,-1)\,\text{per }2_q \\
&\qquad\quad + 32 (1, 0, -1)\,\text{per }3_q \\
&\qquad\quad + 36 (1, 0, -1, 0)\,\text{per }4_q\Bigr\}, \\[1ex]
P(1, 2, 3, 4, 5)q &= \tfrac{1}{86400}\Bigl\{30 q^4 + 900 q^3 + 9300 q^2 + 38250 q + 50651 \\
&\qquad\quad + (1350 q + 10125)(1, -1)\,\text{per }2_q \\
&\qquad\quad + 3200 (2, -1, -1)\,\text{per }3_q \\
&\qquad\quad + 5400 (1, -1, -1, 1)\,\text{per }4_q \\
&\qquad\quad + 3456 (4, -1, -1, -1, -1)\,\text{per }5_q\Bigr\}, \\[1ex]
P(2)q &= \tfrac{1}{2}\Bigl\{1 + (1, -1)\,\text{per }2_q\Bigr\}, \\[1ex]
P(2, 3)q &= \tfrac{1}{12}\Bigl\{2q + 5 + 3 (1, -1)\,\text{per }2_q + 4 (1, -1, 0)\,\text{per }3_q\Bigr\},
\end{align*}
\begin{align*}
P(2, 3, 4)q &= \tfrac{1}{288}\Bigl\{6 q^2 + 54 q + 107 + (18 q + 81)(1, -1)\,\text{per }2_q \\
&\qquad\quad + 32 (2, -1, -1)\,\text{per }3_q + 36 (1, -1, -1, 1)\,\text{per }4_q\Bigr\}, \\[1ex]
P(2, 3, 4, 5)q &= \tfrac{1}{1440}\Bigl\{2 q^3 + 42 q^2 + 267 q + 497 + (45 q + 315)(1, -1)\,\text{per }2_q \\
&\qquad\quad + 160 (1, -1, 0)\,\text{per }3_q + 180 (1, 0, -1, 0)\,\text{per }4_q \\
&\qquad\quad + 288 (1, -1, 0, 0, 0)\,\text{per }5_q\Bigr\}, \\[1ex]
P(2, 3, 4, 5, 6)q &= \tfrac{1}{172800}\Bigl\{10 q^4 + 400 q^3 + 5550 q^2 + 31000 q + 56877 \\
&\qquad\quad + (450 q^2 + 9000 q + 39075)(1, -1)\,\text{per }2_q \\
&\qquad\quad + 3200 q (1, -1, 0)\,\text{per }3_q + 1600 (21, -19, -2)\,\text{per }3_q \\
&\qquad\quad + 10800 (1, 0, -1, 0)\,\text{per }4_q + 6912 (4, -1, -1, -1, -1)\,\text{per }5_q \\
&\qquad\quad + 4800 (1, -1, -2, -1, 1, 2)\,\text{per }6_q\Bigr\}, \\[1ex]
P(1, 2, 3, 5)q &= \tfrac{1}{720}\Bigl\{4 q^3 + 66 q^2 + 324 q + 451 + 45 (1, -1)\,\text{per }2_q \\
&\qquad\quad + 80 (1, -1, 0)\,\text{per }3_q + 144 (1, 0, 0, 0, -1)\,\text{per }5_q\Bigr\}, \\[1ex]
P(1, 2, 2, 3, 4)q &= \tfrac{1}{6912}\Bigl\{6 q^4 + 144 q^3 + 1194 q^2 + 3960 q + 4267 \\
&\qquad\quad + (54 q^2 + 648 q + 1701)(1, -1)\,\text{per }2_q \\
&\qquad\quad + 256 (2, -1, -1)\,\text{per }3_q + 432 (1, 0, -1, 0)\,\text{per }4_q\Bigr\}, \\[1ex]
P(8)q &= \tfrac{1}{8}\Bigl\{1 + 1 (1, -1)\,\text{per }2_q + 2 (1, 0, -1, 0)\,\text{per }4_q \\
&\qquad\quad + 8 (1, 0, 0, 0, -1, 0, 0, 0)\,\text{per }8_q\Bigr\}, \\[1ex]
P(7, 8)q &= \tfrac{1}{112}\Bigl\{2 q + 43 + 7 (1, -1)\,\text{per }2_q + 14 (1, -1, -1, 1)\,\text{per }4_q \\
&\qquad\quad + 16 (3, 2, 1, 0, -1, -2, -3)\,\text{per }7_q \\
&\qquad\quad + 56 (0, -1, -1, 0, 0, 1, 1, 0)\,\text{per }8_q\Bigr\},
\end{align*}
which are, I think, worth preserving.

\bigskip
\noindent\textit{Received April 14,---Read May 3 and 10, 1855.}

\medskip

\noindent I proceed to discuss the following problem: ``To find in how many ways a number $q$ can be made up as a sum of $m$ terms with the elements $0, 1, 2, \ldots k$, each element being repeatable an indefinite number of times.'' The required number of partitions is represented by
\[
P(0, 1, 2, \ldots k)^m q,
\]
and the number of partitions of $q$ less the number of partitions of $q-1$ is represented by
\[
P'(0, 1, 2, \ldots k)^m q.
\]
We have, as is well known,
\[
P(0, 1, 2, \ldots k)^m q = \text{coefficient } x^q z^m \text{ in } \frac{1}{(1-z)(1-xz)\cdots(1-x^k z)},
\]
where the expansion is to be effected in ascending powers of $z$. Now
\[
\frac{1}{(1-z)(1-xz)\cdots(1-x^k z)} = 1 + \frac{1-x^{k+1}}{1-x}z + \frac{(1-x^{k+1})(1-x^{k+2})}{(1-x)(1-x^2)}z^2 + \&\text{c.},
\]
the general term being
\[
\frac{(1-x^{k+1})(1-x^{k+2})\cdots(1-x^{k+s})}{(1-x)(1-x^2)\cdots(1-x^s)}z^s,
\]
or, what is the same thing,
\[
\frac{(1-x^{m+1})(1-x^{m+2})\cdots(1-x^{m+k})}{(1-x)(1-x^2)\cdots(1-x^k)}z^m,
\]
and consequently
\[
P(0, 1, 2, \ldots k)^m q = \text{coefficient } x^q \text{ in } \frac{(1-x^{m+1})(1-x^{m+2})\cdots(1-x^{m+k})}{(1-x)(1-x^2)\cdots(1-x^k)};
\]
to transform this expression I make use of the equation
\[
\frac{1}{(1-z)(1-xz)\cdots(1-x^k z)} = 1 + \frac{1-x^{k+1}}{1-x}z + \frac{(1-x^{k+1})(1-x^{k+2})}{(1-x)(1-x^2)}z^2 + \&\text{c.},
\]
where the general term is
\[
\frac{(1-x^{k+1})(1-x^{k+2})\cdots(1-x^{k+s})}{(1-x)(1-x^2)\cdots(1-x^s)}z^s,
\]
and the series is a finite one, the last term being that corresponding to $s = k$, viz.\ $z^k \div x^{\frac{1}{2}k(k+1)}$. Writing $x^m$ for $z$, and substituting the resulting value of
\[
\frac{(1-x^{m+1})(1-x^{m+2})\cdots(1-x^{m+k})}{(1-x)(1-x^2)\cdots(1-x^k)}
\]
in the formula for $P(0, 1, 2, \ldots k)^m q$, we have
\[
P(0, 1, 2, \ldots k)^m q = \Sigma_s \Bigl\{(-)^s\,\text{coefficient } x^q \text{ in } \frac{x^{sm + \frac{1}{2}s(s+1)}}{(1-x)(1-x^2)\cdots(1-x^s)(1-x)(1-x^2)\cdots(1-x^{k-s})}\Bigr\},
\]
where the summation extends from $s = 0$ to $s = k$; but if for any value of $s$ between these limits $sm + \tfrac{1}{2}s(s+1)$ becomes greater than $q$, then it is clear that the summation need only be extended from $s = 0$ to the last preceding value of $s$, or what is the same thing, from $s = 0$ to the greatest value of $s$ for which $q - sm - \tfrac{1}{2}s(s+1)$ is positive or zero.

It is obvious, that if $q > km$, then
\[
P(0, 1, 2, \ldots k)^m q = 0;
\]
and moreover, that if $q \leq km$, then
\[
P(0, 1, 2, \ldots k)^m q = P(0, 1, 2, \ldots k)^m (km - q),
\]
so that we may always suppose $q \nleq \tfrac{1}{2}km$. I write therefore $q = \tfrac{1}{2}(km - \alpha)$ where $\alpha$ is zero or a positive integer not greater than $km$, and is even or odd according as $km$ is even or odd. Substituting this value of $q$ and making a slight change in the form of the result, we have
\[
P(0, 1, 2, \ldots k)^m \tfrac{1}{2}(km - \alpha) = \Sigma_s\Bigl\{(-)^s\,\text{coeff.\ } x^{\frac{1}{2}k m - \frac{1}{2}\alpha m} \text{ in } \frac{x^{\frac{1}{2}s(s+1)}}{(1-x)\cdots(1-x^s)(1-x)\cdots(1-x^{k-s})}\Bigr\},
\]
where the summation extends from $s = 0$ to the greatest value of $s$ for which $(\tfrac{1}{2}k - s)m - \tfrac{1}{2}\alpha - \tfrac{1}{2}s(s+1)$ is positive or zero. But we may, if we please, consider the summation as extending, when $k$ is even, from $s = 0$ to $s = \tfrac{1}{2}k - 1$, and when $k$ is odd, from $s = 0$ to $s = \tfrac{1}{2}(k-1)$; the terms corresponding to values of $s$ greater than the greatest value for which $(\tfrac{1}{2}k - s)m - \tfrac{1}{2}\alpha - \tfrac{1}{2}s(s+1)$ is positive or zero being of course equal to zero. It may be noticed, that the fraction will be a proper one if $\alpha < (k - s)(k - s + 1)$; or substituting for $s$ its greatest value, the fraction will be a proper one for all values of $s$, if, when $k$ is even, $\alpha < \tfrac{1}{2}k(k + 2)$, and when $k$ is odd, $\alpha < \tfrac{1}{4}(k+1)(k+3)$.

We have in a similar manner,
\[
P'(0, 1, 2, \ldots k)^m q = \text{coefficient } x^q z^m \text{ in } \frac{1-x}{(1-z)(1-xz)\cdots(1-x^k z)},
\]
which leads to
\[
P'(0, 1, 2, \ldots k)^m \tfrac{1}{2}(km - \alpha) = \Sigma_s\Bigl\{(-)^s\,\text{coeff.\ } x^{\frac{1}{2}k m - \frac{1}{2}\alpha m} \text{ in } \frac{x^{\frac{1}{2}s(s+1)}}{(1-x)\cdots(1-x^s)(1-x)\cdots(1-x^{k-s-1})}\Bigr\},
\]
where the summation extends, as in the former case, from $s = 0$ to the greatest value of $s$, for which $(\tfrac{1}{2}k - s)m - \tfrac{1}{2}\alpha - \tfrac{1}{2}s(s+1)$ is positive or zero, or, if we please, when $k$ is even, from $s = 0$ to $s = \tfrac{1}{2}k - 1$, and when $s$ is odd, from $s = 0$ to $s = \tfrac{1}{2}(k - 1)$. The condition, in order that the fraction may be a proper one for all values of $s$, is, when $k$ is even, $\alpha + 1 < \tfrac{1}{2}k(k + 2)$, and when $k$ is odd, $\alpha + 1 < \tfrac{1}{4}(k + 1)(k + 3)$.

To transform the preceding expressions, I write when $k$ is odd $x^2$ instead of $x$, and I put for shortness $\theta$ instead of $\tfrac{1}{2}k - s$ or $\tfrac{1}{2}(k - 2s)$, and $\gamma$ instead of $\tfrac{1}{2}\alpha + \tfrac{1}{2}s(s+1)$ or $\alpha + s(s + 1)$; we have to consider an expression of the form
\[
\text{coefficient } x^{\theta m} \text{ in } \frac{x^\gamma}{F x},
\]
where $F x$ is the product of factors of the form $1 - x^a$. Suppose that $\bar{a}$ is the least common multiple of $a$ and $\theta$, then $(1 - x^{\bar{a}}) \div (1 - x^a)$ is an integral function of $x$, equal to $\chi x$ suppose, and $1 \div (1 - x^a) = \chi x \div (1 - x^{\bar{a}})$. Making this change in all the factors of $F x$ which require it (i.e.\ in all the factors except those in which $a$ is a multiple of $\theta$), the general term becomes
\[
\text{coefficient } x^{\theta m} \text{ in } \frac{x^\gamma\,H x}{G x},
\]
where $G x$ is a product of factors of the form $1 - x^{a'}$, in which $a'$ is a multiple of $\theta$, i.e.\ $G x$ is a rational and integral function of $x^\theta$. But in the numerator $x^\gamma H x$ we may reject, as not contributing to the formation of the coefficient of $x^{\theta m}$, all the terms in which the indices are not multiples of $\theta$; the numerator is thus reduced to a rational and integral function of $x^\theta$, and the general term is therefore of the form
\[
\text{coefficient } x^{\theta m} \text{ in } \frac{\lambda x}{k x},
\]
or what is the same thing, of the form
\[
\text{coefficient } x^m \text{ in } \frac{\lambda x}{k x},
\]
where $k x$ is the product of factors of the form $1 - x^a$, and $\lambda x$ is a rational and integral function of $x$. The particular value of the fraction depends on the value of $s$; and uniting the different terms, we have an expression
\[
\text{coefficient } x^m \text{ in } \Sigma_s\bigl((-)^s\,\tfrac{\lambda x}{k x}\bigr),
\]
which is equivalent to
\[
\text{coefficient } x^m \text{ in } \frac{\phi x}{f x},
\]
where $f x$ is a product of factors of the form $1 - x^a$, and $\phi x$ is a rational and integral function of $x$. And it is clear that the fraction will be a proper one when each of the fractions in the original expression is a proper fraction, i.e.\ in the case of $P(0, 1, 2, \ldots k)^m \tfrac{1}{2}(km - \alpha)$, when for $k$ even, $\alpha < \tfrac{1}{2}k(k+2)$, and for $k$ odd, $\alpha < \tfrac{1}{4}(k+1)(k+3)$; and in the case of $P'(0, 1, 2, \ldots k)^m \tfrac{1}{2}(km - \alpha)$, when for $k$ even, $\alpha + 1 < \tfrac{1}{2}k(k + 2)$, and for $k$ odd, $\alpha + 1 < \tfrac{1}{4}(k + 1)(k + 3)$.

We see, therefore, that $P(0, 1, 2, \ldots k)^m \tfrac{1}{2}(km - \alpha)$, and $P'(0, 1, 2, \ldots k)^m \tfrac{1}{2}(km - \alpha)$, are each of them of the form
\[
\text{coefficient } x^m \text{ in } \frac{\phi x}{f x},
\]
where $f x$ is the product of factors of the form $1 - x^a$, and up to certain limiting values of $\alpha$ the fraction is a proper fraction. When the fraction $\dfrac{\phi x}{f x}$ is known, we may therefore obtain by the method employed in the former part of this Memoir, analytical expressions (involving prime circulators) for the functions $P$ and $P'$.

As an example, take
\[
P(0, 1, 2, 3)^m \tfrac{1}{2}m,
\]
which is equal to
\[
\text{coefficient } x^{\frac{3}{2}m} \text{ in } \frac{1}{(1-x^m)(1-x^2)(1-x^3)} - \text{coefficient } x^{\frac{3}{2}m} \text{ in } \frac{x}{(1-x^2)(1-x^3)(1-x^4)}.
\]

The multiplier for the first fraction is
\[
\frac{(1 - x^6)(1 - x^{12})}{(1 - x^2)(1 - x^4)},
\]
which is equal to
\[
1 + x^2 + 2 x^4 + x^6 + 2 x^8 + x^{10} + x^{12}.
\]
Hence, rejecting in the numerator the terms the indices of which are not divisible by 3, the first term becomes
\[
\text{coefficient } x^{\frac{3}{2}m} \text{ in } \frac{1 + 2 x^6 + x^{12}}{(1 - x^6)(1 - x^{12})},
\]
or what is the same thing, the first term is
\[
\text{coefficient } x^m \text{ in } \frac{1 + 2 x^2 + x^4}{(1 - x^2)^2 (1 - x^4)};
\]
and, the second term being
\[
-\,\text{coefficient } x^{\frac{3}{2}m} \text{ in } \frac{x^4}{(1 - x^6)^2 (1 - x^{12})},
\]
we have
\[
P(0, 1, 2, 3)^m \tfrac{1}{2}m = \text{coefficient } x^m \text{ in } \frac{1 + x^4}{(1 - x^2)^2 (1 - x^4)};
\]
and similarly it may be shown, that
\[
P'(0, 1, 2, 3)^m \tfrac{1}{2}(3m - 1) = \text{coefficient } x^m \text{ in } \frac{1 + x^2}{(1 - x^2)^2(1 - x^4)}.
\]

As another example, take $P'(0, 1, 2, 3, 4, 5)^m \tfrac{1}{2}m$, which is equal to
\begin{align*}
&\text{coefficient } x^{\frac{5}{2}m} \text{ in } \frac{1}{(1-x^m)(1-x^2)(1-x^3)(1-x^4)(1-x^5)} \\
&\quad - \text{coefficient } x^{\frac{5}{2}m} \text{ in } \frac{x}{(1-x^2)(1-x^3)(1-x^4)(1-x^5)(1-x^6)} \\
&\quad + \text{coefficient } x^{\frac{5}{2}m} \text{ in } \frac{x^3}{(1-x^2)(1-x^3)(1-x^4)(1-x^5)(1-x^6)(1-x^7)}.
\end{align*}
The multiplier for the first fraction is
\[
\frac{(1 - x^6)(1 - x^7)(1 - x^8)}{(1 - x^2)(1 - x^3)(1 - x^4)},
\]
which is a function of $x$ of the order 36, the coefficients of which are
\[
1, 0, 1, 1, 2, 1, 3, 2, 4, 3, 4, 4, 6, 4, 6, 5, 7, 5, 7, 5, 7, 5, 6, 4, 6, 4, 4, 3, 4, 2, 3, 1, 2, 1, 1, 0, 1,
\]
and the first part becomes therefore
\[
\text{coefficient } x^{\frac{5}{2}m} \text{ in } \frac{1 + 5 x^{10} + x^{15} + 4 x^{20} + 7 x^{25} + 4 x^{30} + 3 x^{35}}{(1 - x^2)(1 - x^4)(1 - x^6)(1 - x^{10})}.
\]

The multiplier for the second fraction is
\[
\frac{(1 - x^6)(1 - x^7)(1 - x^8)}{(1 - x^3)(1 - x^4)(1 - x^5)},
\]
which is a function of $x$ of the order 14, the coefficients of which are
\[
1, 1, 2, 1, 3, 2, 3, 1, 3, 2, 3, 1, 2, 1, 1;
\]
and the second term becomes
\[
-\,\text{coefficient } x^{\frac{5}{2}m} \text{ in } \frac{x + 2 x^3 + 2 x^5 + 3 x^9 + x^{11} + x^{13}}{(1 - x^2)^2(1 - x^5)(1 - x^{12})},
\]
and the third term is
\[
\text{coefficient } x^{\frac{5}{2}m} \text{ in } \frac{x^3}{(1 - x^4)(1 - x^6)^2}.
\]

Now the fractions may be reduced to a common denominator
\[
(1 - x^2)(1 - x^4)(1 - x^{10})(1 - x^{12}),
\]
by multiplying the terms of the second fraction by $\dfrac{1 - x^4}{1 - x^2}\,(= 1 + x^2 + x^4)$, and the terms of the third fraction by $\dfrac{1 - x^{12}}{1 - x^4}\,(= 1 + x^4)$; performing the operations and adding, the numerator and denominator of the resulting fraction will each of them contain the factor $1 - x^2$; and casting this out, we find
\[
P'(0, 1, 2, 3, 4, 5)^m \tfrac{1}{2}m = \text{coefficient } x^m \text{ in } \frac{x^4 \,(\ldots)}{(1 - x)(1 - x^2)^2(1 - x^4)}.
\]

I have calculated by this method several other particular cases, which are given in my ``Second Memoir upon Quantics'', [141], the present researches were in fact made for the sake of their application to that theory.

\bigskip
\noindent\textit{Received April 20,---Read May 3 and 10, 1855.}

\medskip

\noindent Since the preceding portions of the present Memoir were written, Mr Sylvester has communicated to me a remarkable theorem which has led me to the following additional investigations\footnote{Mr Sylvester's researches are published in the \textit{Quarterly Mathematical Journal}, July 1855, [vol.~i.\ pp.\ 141--152], and he has there given the general formula as well for the circulating as the non-circulating part of the expression for the number of partitions.---Added 23rd February, 1856.---A.~C.}.

Let $\dfrac{\phi x}{f x}$ be a rational fraction, and let $(x - x_1)^k$ be a factor of the denominator $f x$, then if
\[
\left\{\frac{\phi x}{f x}\right\}_{x_1}
\]
denote the portion which is made up of the simple fractions having powers of $x - x_1$ for their denominators, we have by a known theorem
\[
\left\{\frac{\phi x}{f x}\right\}_{x_1} = \Sigma\,\frac{1}{z}\,\text{coefficient } \frac{1}{z} \text{ in } \frac{\phi(x_1 + z)}{x - x_1 - z\,f(x_1 + z)}.
\]
Now by a theorem of Jacobi's and Cauchy's,
\[
\text{coefficient } \frac{1}{z} \text{ in } F z = \text{coefficient } \frac{1}{t} \text{ in } F(\phi^{-1} t)\,\phi' t,
\]
whence, writing $x_1 + z = x_1 e^t$, we have
\[
\left\{\frac{\phi x}{f x}\right\}_{x_1} = \text{coefficient } \frac{1}{t} \text{ in } \frac{x_1 \phi(x_1 e^t) e^t}{x - x_1 e^t\,f(x_1 e^t)}.
\]
Now putting for a moment $x = x_1 e^\theta$, we have
\[
\frac{x_1 - x}{x_1 - x e^t} = \frac{x_1(1 - e^{\theta + t})}{x_1(1 - e^\theta)} = \frac{\theta + t}{\theta}\,\frac{e^{\theta + t}}{e^\theta}\,\text{etc.}
\]
and $e^\theta\,d\theta = x\,dx$, whence
\[
\frac{x_1}{x_1 - x e^t} = \frac{x_1}{x_1 - x}\,\frac{1}{1}\,\frac{x_1 - x}{x_1 - x}\,\frac{1}{1.2}\cdots\,\frac{x_1 - x}{x_1 - x},
\]
the general term of which is
\[
\frac{1}{\Pi(s - 1)}t^{s-1}(x \partial_x)^{s-1}\,\frac{x_1}{x_1 - x}.
\]

Hence representing the general term of $x_1\phi(x_1 e^{-t})$ by $x_1\phi^{(s-1)}\,t^{s-1}$, and writing $\dfrac{x_1\phi(x_1 e^{-t})}{f x} = \gamma x$,
\[
\text{coefficient } \frac{1}{t} \text{ in } \frac{x_1\phi(x_1 e^{-t})}{f x},
\]
we find, writing down only the general term,
\[
\left\{\frac{\phi x}{f x}\right\}_{x_1} = \frac{1}{\Pi(s - 1)}\gamma x^{s-1}\,(x \partial_x)^{s-1}\,\frac{x_1}{x_1 - x},
\]
where the value of $\gamma x$ depends upon that of $s$, and where $s$ extends from $s = 1$ to $s = k$.

Suppose now that the denominator is made up of factors (the same or different) of the form $1 - x^m$. And let $a$ be any divisor of one or more of the indices $m$, and let $k$ be the number of the indices of which $a$ is a divisor. The denominator contains the divisor $[1 - x^a]^k$, and consequently if $\rho$ be any root of the equation $[1 - x^a] = 0$, the denominator contains the factor $(\rho - x)^k$. Hence writing $\rho$ for $x_1$ and taking the sum with respect to all the roots of the equation $[1 - x^a] = 0$, we find
\[
\left\{\frac{\phi x}{f x}\right\}_{[1 - x^a]} = \frac{1}{\Pi(s - 1)}\Sigma\,(x \partial_x)^{s-1}\,\frac{\rho}{\rho - x}
= \cdots = \frac{1}{\Pi(s - 1)}(x \partial_x)^{s-1}\,\frac{\theta x}{[1 - x^a]} + \cdots,
\]
where
\[
\theta_s \rho = \text{coefficient } \frac{1}{t} \text{ in } t^{s-1}\,\psi\rho.
\]
and as before $s$ extends from $s = 1$ to $s = k$. We have thus the actual value of the function $\theta x$ made use of in the memoir.

A preceding formula gives
\[
\left\{\frac{\phi x}{f x}\right\}_{x_1} = \text{coefficient } \frac{1}{t} \text{ in } \frac{1}{1 - x e^t}\,\frac{\phi(x_1 e^{-t})}{f'(x_1 e^{-t})},
\]
which is a very simple expression for the non-circulating part of the fraction $\dfrac{\phi x}{f x}$.

This is, in fact, Mr Sylvester's theorem above referred to.

%==============================================================
% Paper 141
%==============================================================

\bigskip\bigskip
\begin{center}
{\Large\bfseries 141.}\\[1ex]
{\large\bfseries A SECOND MEMOIR UPON QUANTICS.}
\end{center}

\noindent\textit{[From the Philosophical Transactions of the Royal Society of London, vol.~CXLVI.\ for the year, 1856, pp.~101--126. Received April 14,---Read May 24, 1855.]}

\medskip

\noindent The present memoir is intended as a continuation of my Introductory Memoir upon Quantics, t.~CXLIV.\ (1854), p.~245, and must be read in connexion with it; the paragraphs of the two Memoirs are numbered continuously. The special subject of the present memoir is the theorem referred to in the Postscript to the Introductory Memoir, and the various developments arising thereout in relation to the number and form of the covariants of a binary quantic.

\textbf{25.} I have already spoken of asyzygetic covariants and invariants, and I shall have occasion to speak of irreducible covariants and invariants. Considering in general a function $u$ determined like a covariant or invariant by means of a system of partial differential equations, it will be convenient to explain what is meant by an asyzygetic integral and by an irreducible integral. Attending for greater simplicity only to a single set $(a, b, c, \ldots)$, which in the case of the covariants or invariants of a single function will be as before the coefficients or elements of the function, it is assumed that the system admits of integrals of the form $u = P$, $u = Q$, \&c., or as we may express it, of integrals $P$, $Q$, \&c., where $P$, $Q$, \&c.\ are rational and integral homogeneous functions of the set $(a, b, c, \ldots)$, and moreover that the system is such that $P$, $Q$, \&c.\ being integrals, $\phi(P, Q, \&\text{c.})$ is also an integral. Then considering only the integrals which are rational and integral homogeneous functions of the set $(a, b, c, \ldots)$, integrals $P, Q, R, \ldots$ not connected by any linear equation or syzygy (such as $\lambda P + \mu Q + \nu R + \cdots = 0$),\footnote{It is hardly necessary to remark, that the multipliers $\lambda$, $\mu$, $\nu$, \ldots, and generally any coefficients or quantities not expressly stated to contain the set $(a, b, c, \ldots)$, are considered as independent of the set, or to use a convenient word, are considered as ``trivials.''} are said to be asyzygetic; but in speaking of the asyzygetic integrals of a particular degree, it is implied that the integrals are a system such that every other integral of the same degree can be expressed as a linear function (such as $\lambda P + \mu Q + \nu R + \cdots$) of these integrals; and any integral $F$ not expressible as a rational and integral homogeneous function of integrals of inferior degrees is said to be an irreducible integral.

\textbf{26.} Suppose now that $A_1, A_2, A_3$, \&c.\ denote the number of asyzygetic integrals of the degrees $1, 2, 3$, \&c.\ respectively, and let $\alpha_1, \alpha_2, \alpha_3$, \&c.\ be determined by the equations
\begin{align*}
A_1 &= \alpha_1, \\
A_2 &= \tfrac{1}{2}\alpha_1(\alpha_1 + 1) + \alpha_2, \\
A_3 &= \tfrac{1}{6}\alpha_1(\alpha_1 + 1)(\alpha_1 + 2) + \alpha_1 \alpha_2 + \alpha_3, \\
A_4 &= \tfrac{1}{24}\alpha_1(\alpha_1 + 1)(\alpha_1 + 2)(\alpha_1 + 3) + \tfrac{1}{2}\alpha_1(\alpha_1 + 1)\alpha_2 + \alpha_1 \alpha_3 + \tfrac{1}{2}\alpha_2(\alpha_2 + 1) + \alpha_4, \&\text{c.},
\end{align*}
or what is the same thing, suppose that
\[
1 + A_1 x + A_2 x^2 + \&\text{c.} = (1 - x)^{-\alpha_1}(1 - x^2)^{-\alpha_2}(1 - x^3)^{-\alpha_3}\cdots;
\]
a little consideration will show that $\alpha_r$ represents the number of irreducible integrals of the degree $r$ less the number of linear relations or syzygies between the composite or non-irreducible integrals of the same degree. In fact the asyzygetic integrals of the degree 1 are necessarily irreducible, i.e.\ $A_1 = \alpha_1$. Represent for a moment the irreducible integrals of the degree 1 by $X$, $X'$, \&c., then the composite integrals $X^2$, $XX'$, \&c., the number of which is $\tfrac{1}{2}\alpha_1(\alpha_1 + 1)$, must be included among the asyzygetic integrals of the degree 2; and if the composite integrals in question were asyzygetic, there would remain $A_2 - \tfrac{1}{2}\alpha_1(\alpha_1 + 1)$ for the number of irreducible integrals of the degree 2; but if there exist syzygies between the composite integrals in question, the number to be subtracted from $A_2$ will be $\tfrac{1}{2}\alpha_1(\alpha_1 + 1)$ less the number of these syzygies, and we shall have $A_2 - \tfrac{1}{2}\alpha_1(\alpha_1 + 1)$, i.e.\ $\alpha_2$ equal to the number of the irreducible integrals of the degree 2 less the number of syzygies between the composite integrals of the same degree. Again, suppose that $\alpha_2$ is negative $= -\beta_2$, we may for simplicity suppose that there are no irreducible integrals of the degree 2, but that the composite integrals of this degree, $X^2$, $XX'$, \&c., are connected by $\beta_2$ syzygies, such as $\lambda X^2 + \mu XX' + \&\text{c.} = 0$, $\lambda_1 X^2 + \mu_1 XX' + \&\text{c.} = 0$. The asyzygetic integrals of the degree 4 include $X^4$, $X^3 X'$, \&c., the number of which is $\tfrac{1}{24}\alpha_1(\alpha_1 + 1)(\alpha_1 + 2)(\alpha_1 + 3)$; but these composite integrals are not asyzygetic, they are connected by syzygies which are augmentatives of the $\beta_2$ syzygies of the second degree, viz.\ by syzygies such as
\[
(\lambda X^2 + \mu XX' \ldots)X^2 = 0,\ (\lambda X^2 + \mu XX' \ldots)XX' = 0,\ \&\text{c.}\ (\lambda X^2 + \mu XX' \ldots)X^2 = 0,
\]
\[
(\lambda_1 X^2 + \mu_1 XX' \ldots) XX' = 0,\ \&\text{c.},
\]
the number of which is $\tfrac{1}{2}\alpha_1(\alpha_1 + 1)\beta_2$. And these syzygies are themselves not asyzygetic, they are connected by secondary syzygies such as
\[
\lambda(\lambda X^2 + \mu XX' \ldots)X^2 + \mu(\lambda X^2 + \mu XX' \ldots)XX' + \&\text{c.}
\]
\[
- \lambda(\lambda_1 X^2 + \mu_1 XX' \ldots)X^2 - \mu(\lambda_1 X^2 + \mu_1 XX' \ldots)XX' - \&\text{c.} = 0,\ \&\text{c.},
\]
the number of which is $\tfrac{1}{2}\beta_2(\beta_2 - 1)$. The real number of syzygies between the composite integrals $X^4$, $X^3 X'$, \&c.\ (i.e.\ of the syzygies arising out of the $\beta_2$ syzygies between $X^2$, $XX'$, \&c.), is therefore $\tfrac{1}{2}\alpha_1(\alpha_1 + 1)\beta_2 - \tfrac{1}{2}\beta_2(\beta_2 - 1)$, and the number of integrals of the degree 4, arising out of the integrals and syzygies of the degrees 1 and 2 respectively, is therefore
\[
\tfrac{1}{24}\alpha_1(\alpha_1 + 1)(\alpha_1 + 2)(\alpha_1 + 3) + \tfrac{1}{2}\alpha_1(\alpha_1 + 1)\beta_2 + \tfrac{1}{2}\beta_2(\beta_2 - 1)
\]
or writing $-\alpha_2$ instead of $\beta_2$, the number in question is
\[
\tfrac{1}{24}\alpha_1(\alpha_1 + 1)(\alpha_1 + 2)(\alpha_1 + 3) + \tfrac{1}{2}\alpha_1(\alpha_1 + 1)\alpha_2 + \tfrac{1}{2}\alpha_2(\alpha_2 + 1).
\]
The integrals of the degrees 1 and 3 give rise to $\alpha_1\alpha_3$ integrals of the degree 4; and if all the composite integrals obtained as above were asyzygetic, we should have
\[
A_4 = \tfrac{1}{24}\alpha_1(\alpha_1+1)(\alpha_1+2)(\alpha_1+3) + \tfrac{1}{2}\alpha_1(\alpha_1+1)\alpha_2 + \tfrac{1}{2}\alpha_2(\alpha_2+1) + \alpha_1\alpha_3,
\]
i.e.\ $\alpha_4$ as the number of irreducible integrals of the degree 4; but if there exist any further syzygies between the composite integrals, then $\alpha_4$ will be the number of the irreducible integrals of the degree 4 less the number of such further syzygies, and the like reasoning is in all cases applicable.

\textbf{27.} It may be remarked, that for any given partial differential equation, or system of such equations, there will be always a finite number $\nu$ such that given $\nu$ independent integrals every other integral is a function (in general an irrational function only expressible as the root of an equation) of the $\nu$ independent integrals; and if to these integrals we join a single other integral not a rational function of the $\nu$ integrals, it is easy to see that every other integral will be a rational function of the $\nu+1$ integrals; but every such other integral will not in general be a rational and integral function of the $\nu+1$ integrals; and [incorrect] there is not in general any finite number whatever of integrals, such that every other integral is a rational and integral function of these integrals, i.e.\ the number of irreducible integrals is in general infinite; and it would seem that this is in fact the case in the theory of covariants.

\textbf{28.} In the case of the covariants, or the invariants of a binary quantic, $A_n$ is given (this will appear in the sequel) as the coefficient of $x^n$ in the development, in ascending powers of $x$, of a rational fraction $\dfrac{\phi x}{f x}$, where $f x$ is of the form
\[
(1 - x)^{\beta_1}(1 - x^2)^{\beta_2}\cdots(1 - x^k)^{\beta_k},
\]
and the degree of $\phi x$ is less than that of $f x$. We have therefore
\[
1 + A_1 x + A_2 x^2 + \cdots = \frac{\phi x}{f x},
\]
and consequently
\[
\frac{\phi x}{f x} = (1 - x)^{\beta_1 - \alpha_1}(1 - x^2)^{\beta_2 - \alpha_2}\cdots(1 - x^k)^{\beta_k - \alpha_k}(1 - x^{k+1})^{-\alpha_{k+1}}\cdots.
\]

Now every rational factor of a binomial $1 - x^m$ is the irreducible factor of $1 - x^{m'}$, where $m'$ is equal to or a submultiple of $m$. Hence in order that the series $\alpha_1, \alpha_2, \alpha_3, \ldots$ may terminate, $\phi x$ must be made up of factors each of which is the irreducible factor of a binomial $1 - x^m$, or if $\phi x$ be itself irreducible, then $\phi x$ must be the irreducible factor of a binomial $1 - x^m$. Conversely, if $\phi x$ be not of the form in question, the series $\alpha_1, \alpha_2, \alpha_3$, \&c.\ will go on \emph{ad infinitum}, and it is easy to see that there is no point in the series such that the terms beyond that point are all of them negative, i.e.\ there will be irreducible covariants or invariants of indefinitely high degrees; and the number of covariants or invariants will be infinite. The number of invariants is first infinite in the case of a quantic of the seventh order, or septimic; the number of covariants is first infinite in the case of a quantic of the fifth order, or quintic. [As is now well known, these conclusions are incorrect, the number of irreducible covariants or invariants is in every case finite.]

\textbf{29.} Resuming the theory of binary quantics, I consider the quantic
\[
(a, b, \ldots, b', a' \mid x, y)^m.
\]
Here writing
\begin{align*}
[y\partial_x] &= a\partial_b + 2 b\partial_c + \cdots + m b'\partial_{a'} = X, \\
[x\partial_y] &= m b\partial_a + (m-1)c\partial_b + \cdots + a'\partial_{b'} = Y,
\end{align*}
any function which is reduced to zero by each of the operations $X - y\partial_x$, $Y - x\partial_y$ is a covariant of the quantic. But a covariant will always be considered as a rational and integral function separately homogeneous in regard to the facients $(x, y)$ and to the coefficients $(a, b, \ldots, b', a')$. And the words order and degree will be taken to refer to the facients and to the coefficients respectively.

I commence by proving the theorem enunciated, No.~23. It follows at once from the definition, that the covariant is reduced to zero by the operation
\[
X\cdot Y - Y\cdot X - y\partial_x\cdot Y - x\partial_y\cdot X + x\partial_y\cdot y\partial_x + y\partial_x\cdot x\partial_y,
\]
which is equivalent to
\[
X\cdot Y - Y\cdot X + y\partial_y - x\partial_x.
\]
Now
\begin{align*}
X\cdot Y &= XY + X(Y), \\
Y\cdot X &= YX + Y(X),
\end{align*}
where $XY$ and $YX$ are equivalent operations, and
\begin{align*}
X(Y) &= 1\cdot m a\partial_a + 2(m - 1) b\partial_b + \cdots + m\cdot 1\cdot b'\partial_{b'}, \\
Y(X) &= m\cdot 1\cdot b\partial_b + \cdots + 2(m - 1) b'\partial_{b'} + 1\cdot m a'\partial_{a'},
\end{align*}
whence
\[
X(Y) - Y(X) = m a\partial_a + (m - 2)b\partial_b + \cdots - (m - 2)b'\partial_{b'} - m a'\partial_{a'} = k\text{ suppose,}
\]
and the covariant is therefore reduced to zero by the operation
\[
k + y\partial_y - x\partial_x.
\]
Now as regards a term $a^\alpha b^\beta \cdots b'^{\beta'} a'^{\alpha'}\cdot x^i y^j$, we have
\begin{align*}
k &= m\alpha + (m - 2)\beta + \cdots - (m - 2)\beta' - m\alpha', \\
y\partial_y - x\partial_x &= j - i;
\end{align*}
and we see at once that for each term of the covariant we must have
\[
m\alpha + (m - 2)\beta + \cdots - (m - 2)\beta' - m\alpha' + j - i = 0,
\]
i.e.\ if $(x, y)$ are considered as being of the weights $-i, i$ respectively, and $(a, b, \ldots, b', a')$ as being of the weights $-\tfrac{1}{2}m, -\tfrac{1}{2}m + 1, \ldots, \tfrac{1}{2}m - 1, \tfrac{1}{2}m$ respectively, then the weight of each term of the covariant is zero.

But if $(x, y)$ are considered as being of the weights $1, 0$ respectively, and $(a, b, \ldots, b', a')$ as being of the weights $0, 1, \ldots, m-1, m$ respectively, then writing the equation under the form
\[
m(\alpha + \beta + \cdots + \beta' + \alpha') + j + i - 2(0 + \beta + \cdots + (m-1)\beta' + m\alpha' + i) = 0,
\]
and supposing that the covariant is of the order $\mu$ and of the degree $\theta$, each term of the covariant will be of the weight $\tfrac{1}{2}(m\theta + \mu)$.

I shall in the sequel consider the weight as reckoned in the last-mentioned manner. It is convenient to remark, that as regards any function of the coefficients of the degree $\theta$ and of the weight $q$, we have
\[
X\cdot Y - Y\cdot X = m\theta - 2q.
\]

\textbf{30.} Consider now a covariant
\[
(A, B, \ldots, B', A' \mid x, y)^\mu
\]
of the order $\mu$ and of the degree $\theta$; the covariant is reduced to zero by each of the operations $X - y\partial_x$, $Y - x\partial_y$, and we are thus led to the systems of equations
\begin{align*}
XA &= 0, \\
XB &= \mu A, \\
XC &= (\mu - 1)B, \\
&\vdots \\
XB' &= 2 C', \\
XA' &= B'
\end{align*}
and
\begin{align*}
YA &= B, \\
YB &= 2 C, \\
&\vdots \\
YC' &= (\mu - 1)B', \\
YB' &= \mu A', \\
YA' &= 0.
\end{align*}
Conversely if these equations are satisfied the function will be a covariant.

I assume that $A$ is a function of the degree $\theta$ and of the weight $-\tfrac{1}{2}(m\theta - \mu)$, satisfying the condition
\[
XA = 0;
\]
and I represent by $YA, Y^2A, Y^3A$, \&c.\ the results obtained by successive operations with $Y$ upon the function $A$. The function $Y^s A$ will be of the degree $\theta$ and of the weight $\tfrac{1}{2}(m\theta - \mu) + s$. And it is clear that in the series of terms $YA, Y^2A, Y^3A$, \&c., we must at last come to a term which is equal to zero. In fact, since $m$ is the greatest weight of any coefficient, the weight of $Y^s$ is at most equal to $m\theta$, and therefore if $\tfrac{1}{2}(m\theta - \mu) + s > m\theta$, or $s > \tfrac{1}{2}(m\theta + \mu)$, we must have $Y^s A = 0$.

Now writing for greater simplicity $XY$ instead of $X\cdot Y$, and so in similar cases, we have, as regards $Y^s A$,
\[
XY - YX = \mu - 2s.
\]
Hence $(XY - YX)A = \mu A$, and consequently $XYA = YXA + \mu A = \mu A$.
Similarly $(XY - YX) YA = (\mu - 2) YA$, and therefore
\[
XY^2 A = YXYA + (\mu - 2) YA = \mu YA + (\mu - 2)YA = 2(\mu - 1) YA.
\]
And again, $(XY - YX) Y^2 A = (\mu - 4) Y^2 A$, and therefore
\[
XY^3 A = YXY^2A + (\mu - 4) Y^2A = 2(\mu - 1)Y^2A + (\mu - 4)Y^2A = 3(\mu - 2)Y^2A,
\]
or generally
\[
XY^s A = s(\mu - s + 1)Y^{s-1} A.
\]
Hence putting $s = \mu+1, \mu+2, \&\text{c.}$, we have
\begin{align*}
XY^{\mu+1}A &= 0, \\
XY^{\mu+2}A &= (\mu + 2)\cdot 1\cdot Y^{\mu+1}A, \\
XY^{\mu+3}A &= (\mu + 3)\cdot 2\cdot Y^{\mu+2}A,
\end{align*}
equations which show that
\[
Y^{\mu+1}A = 0;
\]
for unless this be so, i.e.\ if $Y^{\mu+1}A \neq 0$, then from the second equation $XY^{\mu+2}A \neq 0$, and therefore $Y^{\mu+2}A \neq 0$, from the third equation $XY^{\mu+3}A \neq 0$, and therefore $Y^{\mu+3}A \neq 0$, and so on \emph{ad infinitum}, i.e.\ we must have $Y^{\mu+1}A = 0$.

\textbf{31.} The suppositions which have been made as to the function $A$, give therefore the equations
\begin{align*}
XA &= 0, \\
XYA &= \mu A, \\
XY^2 A &= 2(\mu - 1)YA,
\end{align*}
and if we now assume $Y^{\mu+1}A = 0$;
\[
B = YA,\quad C = \tfrac{1}{2}YB,\quad \ldots\quad A' = \tfrac{1}{\mu}YB',
\]
the system becomes
\begin{align*}
XA &= 0, \\
XB &= \mu A, \\
XC &= (\mu - 1)B, \\
&\vdots \\
XA' &= B', \\
YA' &= 0;
\end{align*}
so that the entire system of equations which express that $(A, B, \ldots, B', A' \mid x, y)^\mu$ is a covariant is satisfied; hence

\medskip\noindent\textsc{Theorem.} Given a quantic $(a, b, \ldots, b', a' \mid x, y)^m$; if $A$ be a function of the coefficients of the degree $\theta$ and of the weight $\tfrac{1}{2}(m\theta - \mu)$ satisfying the condition $XA = 0$, and if $B, C, \ldots, B', A'$ are determined by the equations
\[
B = YA,\quad C = \tfrac{1}{2}YB,\quad \ldots\quad A' = \tfrac{1}{\mu}YB',
\]
then will
\[
(A, B, \ldots, B', A' \mid x, y)^\mu
\]
be a covariant.

In particular, a function $A$ of the degree $\theta$ and of the weight $\tfrac{1}{2}m\theta$, satisfying the condition $XA = 0$, will (also satisfy the equation $YA = 0$ and will) be an invariant.

\textbf{32.} I take now for $A$ the most general function of the coefficients, of the degree $\theta$ and of the weight $\tfrac{1}{2}(m\theta - \mu)$; then $XA$ is a function of the degree $\theta$ and of the weight $\tfrac{1}{2}(m\theta - \mu) - 1$, and the arbitrary coefficients in the function $A$ are to be determined so that $XA = 0$. The number of arbitrary coefficients is equal to the number of terms in $A$, and the number of the equations to be satisfied is equal to the number of terms in $XA$; hence the number of the arbitrary coefficients which remains indeterminate is equal to the number of terms in $A$ less the number of terms in $XA$; and since the covariant is completely determined when the leading coefficient is known, the difference in question is equal to the number of the asyzygetic covariants, i.e.\ the number of the asyzygetic covariants of the order $\mu$ and the degree $\theta$ is equal to the number of terms of the degree $\theta$ and weight $\tfrac{1}{2}(m\theta - \mu)$, less the number of terms of the degree $\theta$ and weight $\tfrac{1}{2}(m\theta - \mu) - 1$.

\textbf{33.} I shall now give some instances of the calculation of covariants by the method just explained. It is very convenient for this purpose to commence by forming the literal parts by Arbogast's Method of Derivations: we thus form tables such as the following:

\textit{[Figure: Arbogast tables showing literal parts for $(a, b, c)$ quantics up to degree 4, $(a, b, c, d)$ quantics up to degree 4, and $(a, b, c, d, e)$ quantics up to degree 2; each table arranges products of letters by degree (row) and weight (column).]}

\textbf{34.} Thus in the case of a cubic $(a, b, c, d \mid x, y)^3$, the tables show that there will be a single invariant of the degree 4. Represent this by
\begin{align*}
&A\,a^2 d^2 \\
&\quad + B\,abcd \\
&\quad + C\,ac^3 \\
&\quad + D\,b^3 d \\
&\quad + E\,b^2 c^2,
\end{align*}
which is to be operated upon with $a\partial_b + 2b\partial_c + 3c\partial_d$. This gives
\begin{center}
\begin{tabular}{cccc}
$+ B$ & $+ 2 B$ & $+ 6 A$ & $a^2 cd$ \\
        & $+ 6 C$ &         & $a b^2 d$ \\
$+ 3 D$ & $+ 4 E$ & $+ 3 B$ & $abc^2$ \\
$+ 4 E$ &         & $+ 3 D$ & $b^3 c$ \\
\end{tabular}
\end{center}
i.e.\ $B + 6 A = 0$, $3 D + 2 E = 0$, \&c.; or putting $A = 1$, we find $B = -6$, $C = 4$, $D = 4$, $E = -3$, and the invariant is
\begin{align*}
&a^2 d^2 \\
&\quad - 6\,abcd \\
&\quad + 4\,ac^3 \\
&\quad + 4\,b^3 d \\
&\quad - 3\,b^2 c^2.
\end{align*}

Again, there is a covariant of the order 3 and the degree 3. The coefficient of $x^3$ or leading coefficient is
\begin{align*}
&A\,a^2 d \\
&\quad + B\,abc \\
&\quad + C\,b^3
\end{align*}
which operated upon with $a\partial_b + 2b\partial_c + 3c\partial_d$, gives
\begin{center}
\begin{tabular}{ccc}
$+ B$ & $+ 3 A$ & $a^2 c$ \\
$+ 3 C$ & $+ 2 B$ & $a b^2$ \\
\end{tabular}
\end{center}
i.e.\ $B + 3 A = 0$, $3 C + 2 B = 0$; or putting $A = 1$, we have $B = -3$, $C = 2$, and the leading coefficient is
\begin{align*}
&a^2 d \\
&\quad - 3\,abc \\
&\quad + 2\,b^3.
\end{align*}

The coefficient of $x^2 y$ is found by operating upon this with $\tfrac{1}{3}(3 b\partial_a + 2c\partial_b + d\partial_c)$, this gives
\begin{center}
\begin{tabular}{ccc}
$+ 6$ & $- 3$ & $abd$ \\
      & $- 6$ & $ac^2$ \\
$- 9$ & $+ 12$ & $b^2 c$ \\
\end{tabular}
\end{center}
i.e.\ the required coefficient of $x^2 y$ is
\begin{align*}
&3\,abd \\
&\quad - 6\,ac^2 \\
&\quad + 3\,b^2 c;
\end{align*}
and by operating upon this with $\tfrac{1}{2}(3b\partial_a + 2c\partial_b + d\partial_c)$, we have for the coefficient of $xy^2$,
\begin{center}
\begin{tabular}{ccc}
$+ 3$ & $- 6$ & $acd$ \\
      & $+ 9$ & $b^2 d$ \\
$- 9$ & $+ 6$ & $bc^2$ \\
\end{tabular}
\end{center}
i.e.\ the coefficient of $xy^2$ is
\begin{align*}
&-3\,acd \\
&\quad + 6\,b^2 d \\
&\quad - 3\,bc^2.
\end{align*}

Finally, operating upon this with $\tfrac{1}{3}(3b\partial_a + 2c\partial_b + d\partial_c)$, we have for the coefficient of $y^3$,
\begin{center}
\begin{tabular}{ccc}
       & $- 1$ & $ad^2$ \\
$- 3$ & $+ 8$  & $b c d$ \\
       & $- 2$ &         \\
\end{tabular}
\end{center}
i.e.\ the coefficient of $y^3$ is
\begin{align*}
&-\,ad^2 \\
&\quad + 3\,bcd \\
&\quad - 2\,c^3,
\end{align*}
and the covariant is
\[
\left(
\begin{array}{cccc}
a^2 d + 1 & abd + 3 & acd - 3 & ad^2 - 1 \\
          & ac^2 - 6 &        &          \\
abc - 3   &          & b^2 d + 6 & bcd + 3 \\
          & b^2 c + 3 & bc^2 - 3 &         \\
b^3 + 2   &          &          & c^3 - 2
\end{array}
\mid x, y\right)^3.
\]

[I now write the numerical coefficients after instead of before the literal terms.]

I have worked out the example in detail as a specimen of the most convenient method for the actual calculation of more complicated covariants\footnote{Note added Feb.~7, 1856.---The following method for the calculation of an invariant or of the leading coefficient of a covariant, is easily seen to be identical in principle with that given in the text. Write down all the terms of the weight next inferior to that of the invariant or leading coefficient, and operate on each of these separately with the symbol
\[
\text{ind. } b\cdot\frac{a}{b} + 2\,\text{ind. } c\cdot\frac{b}{c} + \cdots + (m-1)\,\text{ind. } b'\cdot\frac{b'}{a'},
\]
where we are first to multiply by the fraction, rejecting negative powers, and then by the index of the proper letter in the term so obtained. Equating the results to zero, we obtain equations between the terms of the invariant or leading coefficient, and replacing in these equations each term by its numerical coefficient in the invariant or leading coefficient, we have the equations of connexion of these numerical coefficients. Thus, for the discriminant of a cubic, the terms of the next inferior weight are $a^2 cd$, $ab^2 d$, $abc^2$, $b^3 c$, and operating on each of these separately with the symbol
\[
\text{ind. } b\cdot\frac{a}{b} + 2\,\text{ind. } c\cdot\frac{b}{c} + 3\,\text{ind. } d\cdot\frac{c}{d},
\]
we find
\begin{center}
\begin{tabular}{lll}
$abcd$        & $+ 4 abcd$ & $+ 6 a^2 d^2$ \\
              & $+ 6 ac^3$ &              \\
$4 b^3 c$    & $+ 4 b^3 c$ & $+ 4 abcd$ \\
$2 b^3 c$    &            &              \\
              &            & $+ 4 b^3 d$ \\
\end{tabular}
\end{center}
and equating the horizontal lines to zero, and assuming $a^2 d^2 = 1$, we have $a^2 d^2 = 1$, $abcd = -6$, $ac^3 = 4$, $b^3 d = 4$, $b^2 c^2 = -3$, or the value of the discriminant is that given in the text.}.

\textbf{35.} The number of terms of the degree $\theta$ and of the weight $q$ is obviously equal to the number of ways in which $q$ can be made up as a sum of $\theta$ terms with the elements $(0, 1, 2, \ldots, m)$, a number which is equal to the coefficient of $x^q z^\theta$ in the development of
\[
\frac{1}{(1 - z)(1 - xz)(1 - x^2 z)\cdots(1 - x^m z)},
\]
and the number of the asyzygetic covariants of any particular degree for the quantic $(*\mid x, y)^m$ can therefore be determined by means of this development. In the case of a cubic, for example, the function to be developed is
\[
\frac{1}{(1 - z)(1 - xz)(1 - x^2 z)(1 - x^3 z)},
\]
which is equal to
\[
1 + (1 + x + x^2 + x^3)z + (1 + x + 2x^2 + 2x^3 + 2x^4 + x^5 + x^6)z^2 + \&\text{c.},
\]
where the coefficients are given by the following table; on account of the symmetry, the series of coefficients for each power of $z$ is continued only to the middle term or middle of the series.

\textit{[Figure: Triangular table of coefficients of the expansion, with rows labelled (0), (1), \ldots, (6) for successive powers of $z$, and entries showing the coefficients of $x^j$ in the corresponding row, mirrored about the middle.]}

\textbf{36.} and from this, by subtracting from each coefficient the coefficient which immediately precedes it, we form the table

\textit{[Figure: Triangular difference table derived from the previous one, with rows labelled (0), (1), \ldots, (6); entries are the successive differences.]}

The successive lines fix the number and character of the covariants of the degrees $0, 1, 2, 3$, \&c. The line $(0)$, if this were to be interpreted, would show that there is a single covariant of the degree $0$; this covariant is of course merely the absolute constant unity, and may be excluded. The line $(1)$ shows that there is a single covariant of the degree 1, viz.\ a covariant of the order 3; this is the cubic itself, which I represent by $U$. The line $(2)$ shows that there are two asyzygetic covariants of the degree 2, viz.\ one of the order 6, this is merely $U^2$, and one of the order 2, this I represent by $H$. The line $(3)$ shows that there are three asyzygetic covariants of the degree 3, viz.\ one of the order 9, this is $U^3$; one of the order 5, this is $UH$, and one of the order 3, this I represent by $\Phi$. The line $(4)$ shows that there are five asyzygetic covariants of the degree 4, viz.\ one of the order 12, this is $U^4$; one of the order 8, this is $U^2 H$; one of the order 6, this is $H^2$; and one of the order 0, i.e.\ an invariant, this I represent by $\nabla$. The line $(5)$ shows that there are six asyzygetic covariants of the degree 5, viz.\ one of the order 15, this is $U^5$; one of the order 11, this is $U^3 H$; one of the order 9, this is $U^2 \Phi$; one of the order 7, this is $UH^2$; one of the order 5, this is $H\Phi$; and one of the order 3, this is $U\nabla$. The line $(6)$ shows that there are 8 asyzygetic covariants of the degree 6, viz.\ one of the order 18, this is $U^6$; one of the order 14, this is $U^4 H$; one of the order 12, this is $U^3\Phi$; one of the order 10, this is $U^2 H^2$; one of the order 8, this is $U H \Phi$; two of the order 6 (i.e.\ the three covariants $H^3$, $\Phi^2$ and $\nabla U^2$ are not asyzygetic, but are connected by a single linear equation or syzygy), and one of the order 2, this is $\nabla H$. We are thus led to the irreducible covariants $U, H, \Phi, \nabla$ connected by a linear equation or syzygy between $H^3$, $\Phi^2$ and $\nabla U^2$, and this is in fact the complete system of irreducible covariants; $\nabla$ is therefore the only invariant.

\textbf{36.} The asyzygetic covariants are of the form $U^p H^q \nabla^r$, or else of the form $\Phi U^p H^q \nabla^r$; and since $U, H, \nabla$ are of the degrees $1, 2, 4$ respectively, and $\Phi$ is of the degree 3, the number of asyzygetic covariants of the degree $m$ of the first form is equal to the coefficient of $x^m$ in $1 \div (1 - x)(1 - x^2)(1 - x^4)$, and the number of the asyzygetic covariants of the degree $m$ of the second form is equal to the coefficient of $x^m$ in $x^3 \div (1 - x)(1 - x^2)(1 - x^4)$. Hence the total number of asyzygetic covariants is equal to the coefficient of $x^m$ in $(1 + x^3) \div (1 - x)(1 - x^2)(1 - x^4)$, or what is the same thing, in
\[
\frac{1}{(1 - x)(1 - x^2)(1 - x^3)(1 - x^4)},
\]
and conversely, if this expression for the number of the asyzygetic covariants of the degree $m$ were established independently, it would follow that the irreducible invariants were four in number, and of the degrees $1, 2, 3, 4$ respectively, but connected by an equation of the degree 6. As regards the invariants, every invariant is of the form $\nabla^p$, i.e.\ the number of asyzygetic invariants of the degree $m$ is equal to the coefficient of $x^m$ in $\dfrac{1}{1 - x^4}$, and conversely, from this expression it would follow that there was a single irreducible invariant of the degree 4.

\textbf{37.} In the case of a quartic, the function to be developed is:
\[
\frac{1}{(1 - z)(1 - xz)(1 - x^2 z)(1 - x^3 z)(1 - x^4 z)},
\]
and the coefficients are given by the table.

\textit{[Figure: Triangular table of coefficients of the expansion for the quartic, with rows labelled (0), (1), \ldots, (6); entries on each row are the coefficients of $x^j$, mirrored about the middle.]}

and subtracting from each coefficient the coefficient immediately preceding it, we have the table

\textit{[Figure: Triangular difference table derived from the previous one, with rows labelled (0), (1), \ldots, (6).]}

the examination of which will show that we have for the quartic the following irreducible covariants, viz.\ the quartic itself $U$; an invariant of the degree 2, which I represent by $I$; a covariant of the order 4 and of the degree 2, which I represent by $H$; an invariant of the degree 3, which I represent by $J$; and a covariant of the order 6 and the degree 3, which I represent by $\Phi$; but that the irreducible covariants are connected by an equation of the degree 6, viz.\ there is a linear equation or syzygy between $\Phi^2$, $I^2 H^2$, $I J H^2 U$, $J H^2 U^2$ and $J^2 U^2$; this is in fact the complete system of the irreducible covariants of the quartic: the only irreducible invariants are the invariants $I, J$.

\textbf{38.} The asyzygetic covariants are of the form $U^p I^q H^r J^s$, or else of the form $\Phi U^p I^q H^r J^s$, and the number of the asyzygetic covariants of the degree $m$ is equal to the coefficient of $x^m$ in $(1 + x^3) \div (1 - x)(1 - x^2)^2(1 - x^3)$, or what is the same thing, in
\[
\frac{1 - x^6}{(1 - x)(1 - x^2)^2(1 - x^3)^2},
\]
and the asyzygetic invariants are of the form $I^p J^q$, and the number of the asyzygetic invariants of the degree $m$ is equal to the coefficient of $x^m$ in $1 \div (1 - x^2)(1 - x^3)$. Conversely, if these formulae were established, the preceding results as to the form of the system of the irreducible covariants or of the irreducible invariants, would follow.

\textbf{39.} In the case of a quintic, the function to be developed is
\[
\frac{1}{(1 - z)(1 - xz)(1 - x^2 z)(1 - x^3 z)(1 - x^4 z)(1 - x^5 z)},
\]
and the coefficients are given by the table

\textit{[Figure: Triangular table of coefficients of the expansion for the quintic, with rows labelled (0), (1), \ldots, (5).]}

and subtracting from each coefficient the one which immediately precedes it, we have the table

\textit{[Figure: Triangular difference table for the quintic, rows labelled (0), (1), \ldots, (5).]}

We thus obtain the following irreducible covariants, viz.:

Of the degree 1; a single covariant of the order 5, this is the quintic itself.

Of the degree 2; two covariants, viz.\ one of the order 6, and one of the order 2.

Of the degree 3; three covariants, viz.\ one of the order 9, one of the order 5, and one of the order 3.

Of the degree 4; three covariants, viz.\ one of the order 6, one of the order 4, and one of the order 0 (an invariant).

Of the degree 5; three covariants, viz.\ one of the order 7, one of the order 3, and one of the order 1 (a linear covariant).

These covariants are connected by a single syzygy of the degree 5 and of the order 11; in fact, the table shows that there are only two asyzygetic covariants of this degree and order; but we may, with the above-mentioned irreducible covariants of the degrees, 1, 2, 3 and 4, form three covariants of the degree 5 and the order 11; there is therefore a syzygy of this degree and order.

\textbf{40.} I represent the number of ways in which $q$ can be made up as a sum of $m$ terms with the elements $0, 1, 2, \ldots, m$, each element being repeatable an indefinite number of times by the notation
\[
P(0, 1, 2, \ldots, m)^\theta q,
\]
and I write for shortness
\[
P'(0, 1, 2, \ldots, m)^\theta q = P(0, 1, 2, \ldots, m)^\theta q - P(0, 1, 2, \ldots, m)^\theta (q - 1).
\]
Then for a quantic of the order $m$, the number of asyzygetic covariants of the degree $\theta$ and of the order $\mu$ is
\[
P'(0, 1, 2, \ldots, m)^\theta\,\tfrac{1}{2}(m\theta - \mu).
\]
In particular, the number of asyzygetic invariants of the degree $\theta$ is
\[
P'(0, 1, 2, \ldots, m)^\theta\,\tfrac{1}{2}m\theta.
\]
To find the total number of the asyzygetic covariants of the degree $\theta$, suppose first that $m\theta$ is even; then, giving to $\mu$ the successive values $0, 2, 4, \ldots, m\theta$, the required number is
\begin{align*}
&P(\tfrac{1}{2}m\theta) - P(\tfrac{1}{2}m\theta - 1) \\
&\quad + P(\tfrac{1}{2}m\theta - 1) - P(\tfrac{1}{2}m\theta - 2) \\
&\quad \vdots \\
&\quad + P(2) - P(1) \\
&\quad + P(1) \\
&= P(\tfrac{1}{2}m\theta),
\end{align*}
i.e.\ when $m\theta$ is even, the number of the asyzygetic covariants of the degree $\theta$ is
\[
P(0, 1, 2, \ldots, m)^\theta\,\tfrac{1}{2}m\theta;
\]
and similarly, when $m\theta$ is odd, the number of the asyzygetic covariants of the degree $\theta$ is
\[
P(0, 1, 2, \ldots, m)^\theta\,\tfrac{1}{2}(m\theta - 1).
\]
But the two formulae may be united into a single formula; for when $m\theta$ is odd $\tfrac{1}{2}m\theta$ is a fraction, and therefore $P(\tfrac{1}{2}m\theta)$ vanishes, and so when $m\theta$ is even $\tfrac{1}{2}(m\theta - 1)$ is a fraction, and $P\tfrac{1}{2}(m\theta - 1)$ vanishes; we have thus the theorem, that for a quantic of the order $m$:

The number of the asyzygetic covariants of the degree $\theta$ is
\[
P(0, 1, 2, \ldots, m)^\theta\,\tfrac{1}{2}m\theta + P(0, 1, 2, \ldots, m)^\theta\,\tfrac{1}{2}(m\theta - 1).
\]

\textbf{41.} The functions $P(\tfrac{1}{2}m\theta)$, \&c.\ may, by the method explained in my ``Researches on the Partition of Numbers,'' [140], be determined as the coefficients of $x^\theta$ in certain functions of $x$; I have calculated the following particular cases:

Putting, for shortness,
\[
P'(0, 1, 2, \ldots, m)^\theta\,\tfrac{1}{2}m\theta = \text{coefficient } x^\theta \text{ in } \phi m,
\]
then
\begin{align*}
\phi 2 &= \frac{1}{1 - x^2}, \\
\phi 3 &= \frac{1}{1 - x^4}, \\
\phi 4 &= \frac{1}{(1 - x^2)(1 - x^3)}, \\
\phi 5 &= \frac{1 - x^8 + x^{18}}{(1 - x^4)(1 - x^8)(1 - x^{12})}, \\
\phi 6 &= \frac{(1 - x)(1 + x - x^2 - x^4 - x^5 + x^7 + x^8)}{(1 - x^2)^2(1 - x^3)(1 - x^4)(1 - x^5)}, \\
\phi 7 &= \frac{1 - x^4 + 2 x^8 - x^{10} + 5 x^{12} + 2 x^{14} + 6 x^{16} + 2 x^{18} + 5 x^{20} - x^{22} + 2 x^{24} - x^{28} + x^{32}}{(1 - x^4)(1 - x^8)(1 - x^{12})(1 - x^{14})(1 - x^{18})}, \\
\phi 8 &= \frac{(1 - x)(1 + x - x^2 - x^4 + x^7 + x^8 + x^9 + x^{10} - x^{11} + x^{13} + x^{14})}{(1 - x^2)^2(1 - x^3)^2(1 - x^4)(1 - x^5)(1 - x^6)},
\end{align*}
and putting
\[
P(0, 1, 2, \ldots, m)^\theta\,\tfrac{1}{2}m\theta = \text{coefficient of } x^\theta \text{ in } \psi m,
\]
then
\begin{align*}
\psi 2 &= \frac{1}{(1 - x)(1 - x^2)}, \\
\psi 3 &= \frac{1 + x^3}{(1 - x^2)(1 - x^4)}, \\
\psi 4 &= \frac{1 - x + x^2}{(1 - x)^2(1 - x^2)(1 - x^3)}, \\
\psi 5 &= \frac{1 + x^4 + 6 x^6 + 9 x^7 + 12 x^8 + 9 x^9 + 6 x^{10} + x^{12} + x^{16}}{(1 - x^2)^2(1 - x^4)(1 - x^6)(1 - x^8)},
\end{align*}
and putting
\[
P(0, 1, 2, \ldots, m)^\theta\,\tfrac{1}{2}(m\theta - 1) = \text{coefficient of } x^\theta \text{ in } \chi m,
\]
then
\begin{align*}
\chi 3 &= \frac{x + x^3}{(1 - x^2)^2(1 - x^3)}, \\
\chi 5 &= \frac{x + 4 x^2 + 8 x^3 + 10 x^4 + 10 x^5 + 8 x^6 + 4 x^7 + x^8}{(1 - x^2)^2(1 - x^4)(1 - x^6)(1 - x^8)}.
\end{align*}

And from what has preceded, it appears that for a quantic of the order $m$, the number of asyzygetic covariants of the degree $\theta$ is for $m$ even, coefficient $x^\theta$ in $\psi m$, and for $m$ odd, coefficient $x^\theta$ in $(\psi m + \chi m)$; and that the number of asyzygetic invariants of the degree $\theta$ is coefficient $x^\theta$ in $\phi m$. Attending first to the invariants:

\textbf{42.} For a quadric, the number of asyzygetic invariants of the degree $\theta$ is
\[
\text{coefficient } x^\theta \text{ in } \frac{1}{1 - x^2},
\]
which leads to the conclusion that there is a single irreducible invariant of the degree 2.

\textbf{43.} For a cubic, the number of asyzygetic invariants of the degree $\theta$ is
\[
\text{coefficient } x^\theta \text{ in } \frac{1}{1 - x^4},
\]
i.e.\ there is a single irreducible invariant of the degree 4.

\textbf{44.} For a quartic, the number of asyzygetic invariants of the degree $\theta$ is
\[
\text{coefficient } x^\theta \text{ in } \frac{1}{(1 - x^2)(1 - x^3)},
\]
i.e.\ there are two irreducible invariants of the degrees 2 and 3 respectively.

\textbf{45.} For a quintic, the number of asyzygetic invariants of the degree $\theta$ is
\[
\text{coefficient } x^\theta \text{ in } \frac{1 - x^8 + x^{18}}{(1 - x^4)(1 - x^8)(1 - x^{12})}.
\]
The numerator is the irreducible factor of $1 - x^{36}$, i.e.\ it is equal to $(1 - x^{36})(1 - x^4) \div (1 - x^{12})(1 - x^9)$; and substituting this value, the number becomes
\[
\text{coefficient } x^\theta \text{ in } \frac{1 - x^{36}}{(1 - x^4)(1 - x^8)(1 - x^{12})(1 - x^{18})},
\]
i.e.\ there are in all four irreducible invariants, which are of the degrees $4, 8, 12$ and $18$ respectively; but these are connected by an equation of the degree 36, i.e.\ the square of the invariant of the degree 18 is a rational and integral function of the other three invariants; a result, the discovery of which is due to M.\ Hermite.

\textbf{46.} For a sextic, the number of asyzygetic invariants of the degree $\theta$ is
\[
\text{coefficient } x^\theta \text{ in } \frac{(1 - x)(1 + x - x^2 - x^4 - x^5 + x^7 + x^8)}{(1 - x^2)^2(1 - x^3)(1 - x^4)(1 - x^5)};
\]
the second factor of the numerator is the irreducible factor $1 - x^{30}$, i.e.\ it is equal to $(1 - x^{30})(1 - x^5)(1 - x^3)(1 - x^2) \div (1 - x^{15})(1 - x^{10})(1 - x^6)(1 - x)$; and substituting this value, the number becomes
\[
\text{coefficient } x^\theta \text{ in } \frac{1 - x^{30}}{(1 - x^2)(1 - x^4)(1 - x^6)(1 - x^{10})(1 - x^{15})},
\]
i.e.\ there are in all five irreducible invariants, which are of the degrees $2, 4, 6, 10$ and $15$ respectively; but these are connected by an equation of the degree 30, i.e.\ the square of the invariant of the degree 15 is a rational and integral function of the other four invariants.

\textbf{47.} For a septimic, the number of asyzygetic invariants of the degree $\theta$ is
\[
\text{coefficient } x^\theta \text{ in } \frac{1 - x^4 + 2 x^8 - x^{10} + 5 x^{12} + 2 x^{14} + 6 x^{16} + 2 x^{18} + 5 x^{20} - x^{22} + 2 x^{24} - x^{28} + x^{32}}{(1 - x^4)(1 - x^8)(1 - x^{12})(1 - x^{14})(1 - x^{18})};
\]
the numerator is equal to
\[
(1 - x^4)(1 - x^8)^{-1}(1 - x^{10})(1 - x^{12})^{-1}(1 - x^{14})^{-1}\cdots,
\]
where the series of factors does not terminate; hence [incorrect, see p.~253] the number of irreducible invariants is infinite; substituting the preceding value, the number of asyzygetic invariants of the degree $\theta$ is
\[
\text{coefficient } x^\theta \text{ in } (1 - x^4)^{-1}(1 - x^8)^{-3}(1 - x^{12})^{-6}(1 - x^{14})^{-4}\cdots.
\]
The first four indices give the number of irreducible invariants of the corresponding degrees, i.e.\ there are $1, 3, 6$ and $4$ irreducible invariants of the degrees $4, 8, 12$ and $14$ respectively, but there is no reason to believe that the same thing holds with respect to the indices of the subsequent terms. To verify this it is to be remarked, that there are $1, 4, 10$ and $4$ asyzygetic invariants of the degrees in question respectively; there is therefore one irreducible invariant of the degree 4; calling this $X_4$, there is only one composite invariant of the degree 8, viz.\ $X_4^2$; there are therefore three irreducible invariants of this degree, say $X_8, X_8', X_8''$. The composite invariants of the degree 12 are four in number, viz.\ $X_4^3, X_4 X_8, X_4 X_8', X_4 X_8''$, and these cannot be connected by any syzygy, for if they were so, $X_4^2, X_8, X_8', X_8''$ would be connected by a syzygy, or there would be less than 3 irreducible invariants of the degree 8. Hence there are precisely 6 irreducible invariants of the degree 12. And since the irreducible invariants of the degrees 4, 8 and 12 do not give rise to any composite invariant of the degree 14, there are precisely 4 irreducible invariants of the degree 14.

\textbf{48.} For an octavic, the number of the asyzygetic invariants of the degree $\theta$ is
\[
\text{coefficient } x^\theta \text{ in } \frac{(1 - x)(1 + x - x^2 - x^4 + x^7 + x^8 + x^9 + x^{10} - x^{11} + x^{13} + x^{14})}{(1 - x^2)^2(1 - x^3)^2(1 - x^4)(1 - x^5)(1 - x^6)},
\]
and the second factor of the numerator is
\[
(1 - x)^{-1}(1 - x^2)(1 - x^3)^{-1}(1 - x^4)^{-1}(1 - x^5)^{-1}(1 - x^6)^{-1}(1 - x^{10})^{-1}(1 - x^{12})(1 - x^{14})(1 - x^{15})\cdots,
\]
where the series of factors does not terminate, hence [incorrect] the number of irreducible invariants is infinite. Substituting the preceding value, the number of the asyzygetic invariants of the degree $\theta$ is
\[ \resizebox{\textwidth}{!}{$\displaystyle
\text{coefficient } x^\theta \text{ in } (1 - x^2)^{-1}(1 - x^3)^{-1}(1 - x^4)^{-1}(1 - x^5)^{-1}(1 - x^6)^{-1}(1 - x^7)^{-2}(1 - x^8)^{-3}(1 - x^9)^{-5}(1 - x^{10})^{-7}\cdots.
$} \]
There is certainly one, and only one irreducible invariant for each of the degrees $2, 3, 4, 5$ and $6$ respectively; but the formula does not show the number of the irreducible invariants of the degrees $7, \&\text{c.}$; in fact, representing the irreducible invariants of the degrees $2, 3, 4, 5$ and $6$ by $X_2, X_3, X_4, X_5, X_6$, these give rise to 3 composite invariants of the degree 7, viz.\ $X_2 X_2 X_3, X_3 X_4, X_2 X_5$, which may or may not be connected by a syzygy; if they are not connected by a syzygy, there will be a single irreducible invariant of the degree 7; but if they are connected by a syzygy, there will be two irreducible invariants of the degree 7; it is useless at present to pursue the discussion further.

Considering next the covariants,

\textbf{49.} For a quadric, the number of asyzygetic covariants of the degree $\theta$ is
\[
\text{coefficient } x^\theta \text{ in } \frac{1}{(1 - x)(1 - x^2)},
\]
i.e.\ there are two irreducible covariants of the degrees 1 and 2 respectively; these are of course the quadric itself and the invariant.

\textbf{50.} For a cubic, the number of the asyzygetic covariants of the degree $\theta$ is
\[
\text{coefficient } x^\theta \text{ in } \frac{(1 + x)(1 + x^3)}{(1 - x^2)^2(1 - x^4)}.
\]
The first factor of the numerator is the irreducible factor of $1 - x^2 = (1 - x^2) \div (1 - x)$, and the second factor of the numerator is the irreducible factor of $1 - x^6 = (1 - x^6) \div (1 - x^2)$; substituting these values, the number is
\[
\text{coefficient } x^\theta \text{ in } \frac{1 - x^6}{(1 - x)(1 - x^2)(1 - x^3)(1 - x^4)},
\]
i.e.\ there are 4 irreducible covariants of the degrees $1, 2, 3, 4$ respectively; but these are connected by an equation of the degree 6; the covariant of the degree 1 is the cubic itself $U$, the other covariants are the covariants already spoken of and represented by the letters $H, \Phi$ and $\nabla$ respectively ($H$ is of the degree 2 and the order 3, $\Phi$ of the degree 3 and the order 3, and $\nabla$ is of the degree 4 and the order 0, i.e.\ it is an invariant).

\textbf{51.} For a quartic, the number of the asyzygetic covariants of the degree $\theta$ is
\[
\text{coefficient } x^\theta \text{ in } \frac{1 - x + x^2}{(1 - x)^2(1 - x^2)(1 - x^3)};
\]
the numerator of which is the irreducible factor of $1 - x^6$, i.e.\ it is equal to $(1 - x^6)(1 - x) \div (1 - x^2)(1 - x^3)$. Making this substitution, the number is
\[
\text{coefficient } x^\theta \text{ in } \frac{1 - x^6}{(1 - x)(1 - x^2)^2(1 - x^3)^2},
\]
i.e.\ there are five irreducible covariants, one of the degree 1, two of the degree 2, and two of the degree 3, but these are connected by an equation of the degree 6. The irreducible covariant of the degree 1 is of course the quartic itself $U$, the other irreducible covariants are those already spoken of and represented by $I, H, J, \Phi$ respectively ($I$ is of the degree 2 and the order 0, and $J$ is of the degree 3 and the order 0, i.e.\ $I$ and $J$ are invariants, $H$ is of the degree 2 and the order 4, $\Phi$ is of the degree 3 and the order 6).

\textbf{52.} For a quintic, the number of irreducible covariants of the degree $\theta$ is
\[ \resizebox{\textwidth}{!}{$\displaystyle
\text{coefficient } x^\theta \text{ in } \frac{1 + x + x^2 + 4 x^3 + 6 x^4 + 8 x^5 + 9 x^6 + 10 x^7 + 12 x^8 + 10 x^9 + 9 x^{10} + 8 x^{11} + 6 x^{12} + 4 x^{13} + x^{14} + x^{15} + x^{16}}{(1 - x^2)^2(1 - x^4)(1 - x^6)(1 - x^8)};
$} \]
the numerator of which is
\[
(1 + x)^2(1 + x - x^2 - x^4 + 2 x^5 + 3 x^6 + 5 x^7 + 4 x^8 + 3 x^9 + 2 x^{10} - x^{12} - x^{13} + x^{14}).
\]
The first factor is $(1 - x^2)^2 \div (1 - x)^2$, the second factor is
\[
(1 - x^2)^3(1 - x^3)^2(1 - x^4)^3(1 - x^5)^4(1 - x^7)^{-5}(1 - x^8)^{-7}\cdots,
\]
which does not terminate; hence [incorrect] the number of irreducible covariants is infinite. Substituting the preceding values, the expression for the number of the asyzygetic covariants of the degree $\theta$ is
\[
\text{coeff. } x^\theta \text{ in } (1 - x)^{-1}(1 - x^2)^{-2}(1 - x^3)^{-3}(1 - x^4)^{-3}(1 - x^5)^{-3}(1 - x^6)^{-4}(1 - x^7)^{-5}(1 - x^8)^{-7}\cdots,
\]
which agrees with a previous result: the numbers of irreducible covariants for the degrees $1, 2, 3, 4$ are $1, 2, 3$ and $3$ respectively, and for the degree 5, the number of irreducible covariants is three, but there is one syzygy between the composite covariants of the degree in question; the difference $3 - 1 = 2$ is the index taken with its sign reversed of the factor $(1 - x^5)^{-3}$.

\textbf{53.} I consider a system of the asyzygetic covariants of any particular degree and order of a given quantic, the system may of course be replaced by a system the terms of which are any linear functions of those of the original system, and it is necessary to inquire what covariants ought to be selected as most proper to represent the system of asyzygetic covariants; the following considerations seem to me to furnish a convenient rule of selection. Let the literal parts of the terms which enter into the coefficients of the highest power of $x$ or leading coefficients be represented by $M_\alpha, M_\beta, M_\gamma, \ldots$ these quantities being arranged in the natural or alphabetical order; the first in order of these quantities $M$, which enters into the leading coefficient of a particular covariant, may for shortness be called the leading term of such covariant, and a covariant the leading term of which is posterior in order to the leading term of another covariant, may be said to have a lower leading term.

It is clear, that by properly determining the multipliers of the linear functions we may form a covariant the leading term of which is lower than the leading term of any other covariant (the definition implies that there is but one such covariant); call this $\Theta_1$. We may in like manner form a covariant such that its leading term is lower than the leading term of every other covariant except $\Theta_1$; or rather we may form a system of such covariants, since if $\Theta_2$ be a covariant having the property in question, $\Theta_2 + k\Theta_1$ will have the same property, but $k$ may be determined so that the covariant shall not contain the leading term of $\Theta_1$, i.e.\ we may form a covariant $\Theta_2$ such that its leading term is lower than the leading term of every other covariant excepting $\Theta_1$, and that the leading term of $\Theta_1$ does not enter into $\Theta_2$; and there is but one such covariant, $\Theta_2$. Again, we may form a covariant $\Theta_3$ such that its leading term is lower than the leading term of every other covariant excepting $\Theta_1$ and $\Theta_2$, and that the leading terms of $\Theta_1$ and $\Theta_2$ do not either of them enter into $\Theta_3$; and there is but one such covariant, $\Theta_3$. And so on, until we arrive at a covariant the leading term of which is higher than the leading terms of the other covariants, and which does not contain the leading terms of the other covariants. We have thus a series of covariants $\Theta_1, \Theta_2, \Theta_3, \&\text{c.}$ containing the proper number of terms, and which covariants may be taken to represent the asyzygetic covariants of the degree and order in question.

In order to render the covariants $\Theta$ definite as well numerically as in regard to sign, we may suppose that the covariant is in its least terms (i.e.\ we may reject numerical factors common to all the terms), and we may make the leading term positive. The leading term with the proper numerical coefficient (if different from unity) and with the proper power of $x$ (or the order of the function) annexed, will, when the covariants of a quantic are tabulated, be sufficient to indicate, without any ambiguity whatever, the particular covariant referred to. I subjoin a table of the covariants of a quadric, a cubic and a quartic, and of the covariants of the degrees $1, 2, 3, 4$ and $5$ respectively of a quintic, and also two other invariants of a quintic.

[Except for the quantic itself, the algebraical sum of the numerical coefficients in any column is $=0$, viz.\ the sum of the coefficients with the sign $+$ is equal to that of the coefficients with the sign $-$, and I have as a numerical verification inserted at the foot of each column this sum with the sign $\pm$].

\bigskip
\begin{center}\textbf{Covariant Tables (Nos.\ 1 to 26).}\end{center}

\textbf{No.~1.} \quad $(a+1,\ b+2,\ c+1 \mid x, y)^2$.

\textbf{No.~2.} \quad $\bigl(ac + 1,\ b^2 - 1\bigr) \quad \pm 1$.

The tables Nos.~1 and 2 are the covariants of a binary quadric. No.~1 is the quadric itself; No.~2 is the quadrinvariant, which is also the discriminant.

\textbf{No.~3.} \quad $(a+1,\ b+3,\ c+3,\ d+1 \mid x, y)^3$.

\textbf{No.~4.} \quad
$\left(\begin{array}{ccc} ac+1 & ad+1 & bd+1 \\ & bc-1 & \\ b^2-1 & & c^2-1 \end{array} \mid x, y\right)^2$, \quad $\pm 1,\ \pm 1,\ \pm 1$.

\textbf{No.~5.} \quad
$\left(\begin{array}{cccc} a^2 d + 1 & abd + 3 & acd - 3 & ad^2 - 1 \\ abc - 3 & ac^2 - 6 & b^2 d + 6 & bcd + 3 \\ b^3 + 2 & b^2 c + 3 & bc^2 - 3 & c^3 - 2 \end{array} \mid x, y\right)^3$, \quad $\pm 3,\ \pm 6,\ \pm 6,\ \pm 3$.

\textbf{No.~6.} \quad
\begin{tabular}{l}
$a^2 d^2 + 1$ \\
$abcd - 6$ \\
$ac^3 + 4$ \\
$b^3 d + 4$ \\
$b^2 c^2 - 3$ \\
\hline
$\pm 7$
\end{tabular}

The tables Nos.~3, 4, and 5 are the covariants of a binary cubic. No.~3 is the cubic itself; No.~4 is the quadricovariant, or Hessian; No.~5 is the cubicovariant; No.~6 is the invariant, or discriminant. And if we write No.~3 $= U$, No.~4 $= H$, No.~5 $= \Phi$, No.~6 $= \nabla$, then identically,
\[
\Phi^2 - \nabla U^2 + 4 H^3 = 0.
\]

\textbf{No.~7.} \quad $(a+1,\ b+4,\ c+6,\ d+4,\ e+1 \mid x, y)^4$.

\textbf{No.~8.} \quad
$\left(\begin{array}{c} ae + 1 \\ bd - 4 \\ c^2 + 3 \end{array}\right) \quad \pm 4$.

\textbf{No.~9.} \quad
$\left(\begin{array}{cccc} ac + 1 & ad + 1 & ae + 1 & bd + 1 \\ & & bd + 2 & be + 2 \\ b^2 - 1 & bc - 2 & c^2 - 2 & cd - 2 \\ & & & \\ \pm 1 & \pm 2 & \pm 3 & \pm 2 \end{array}\right)\ +\ (be + 2,\ ce - 1)$.

\textit{[Tables 10, 11, 12 follow with analogous formats: cubinvariant, cubicovariant, and discriminant of the quartic. The discriminant (No.~12) has 18 terms with sum $\pm 442$.]}

The tables Nos.~7, 8, 9, 10 and 11 are the irreducible covariants of a quartic. No.~7 is the quartic itself; No.~8 is the quadrinvariant; No.~9 is the quadricovariant, or Hessian; No.~10 is the cubinvariant; and No.~11 is the cubicovariant. The table No.~12 is the discriminant. And if we write No.~7 $= U$, No.~8 $= I$, No.~9 $= H$, No.~10 $= J$, No.~11 $= \Phi$, No.~12 $= \nabla$, then the irreducible covariants are connected by the identical equation
\[
JU^3 - IU^2 H + 4 H^3 + \Phi^2 = 0,
\]
and we have
\[
\nabla = I^3 - 27 J^2.
\]

[The Tables Nos.~13 to 24 which follow, and also Nos.~25 and 26 which are given in 143 relate to the binary quintic. I have inserted in the headings the capital letters $A, B, \ldots, L$ and also $Q$ and $Q'$ by which I refer to these covariants of the quintic. $A$ is the quintic itself, $C$ is the Hessian, $G$ is the quartinvariant, $J$ a linear covariant; $Q$ is the simplest octinvariant, and $Q'$ is the discriminant. As noticed in the original memoir we have $AI + BF - CE = 0$; and $Q' = G^2 - 128 Q$, only the coefficient 128 was by mistake given as 1152.]

\textbf{A.~No.~13.} \quad $(a+1,\ b+5,\ c+10,\ d+10,\ e+5,\ f+1 \mid x, y)^5$.

\textbf{B.~No.~14.} \quad
$\left(\begin{array}{ccc} ae + 1 & af + 1 & bf + 1 \\ bd - 4 & be - 3 & ce - 4 \\ c^2 + 3 & cd + 2 & d^2 + 3 \\ \pm 4 & \pm 3 & \pm 4 \end{array} \mid x, y\right)^2$.

\textbf{C.~No.~15.} \quad
$\left(\begin{array}{ccccccc}
ac + 1 & ad + 3 & ae + 3 & af + 1 & bf + 3 & cf + 3 & df + 1 \\
        & bd + 3 & bd + 7 & & ce + 3 & de - 3 & \\
b^2 - 1 & bc - 3 & be - 6 & cd - 8 & & & c^2 - 1 \\
        &        & c^2 - 6 & d^2 - 6 & & & \\
\pm 1 & \pm 3 & \pm 6 & \pm 8 & \pm 6 & \pm 3 & \pm 1
\end{array} \mid x, y\right)^6$.

\textit{[Tables D (No.~16), E (No.~17), F (No.~18), G (No.~19), H (No.~20), I (No.~21), J (No.~22), K (No.~23), L (No.~24) follow with the full literal expansions. The detailed numerical tables for these covariants of the binary quintic are extensive; the column sums (with sign $\pm$) shown at the bottom of each column serve as numerical verification. Several of the larger tables involve products of $a$, $b$, \ldots, $f$ of degree up to 5 and weight up to 16.]}

\textbf{No.~25.} \quad $Q = + 1\,a^2 cdf + \&\text{c.}$ \quad [see post, 143.]

\textbf{No.~26.} \quad $Q' = + 1\,a^4 f^4 + \&\text{c.}$

%==============================================================
% Paper 142
%==============================================================

\bigskip\bigskip
\begin{center}
{\Large\bfseries 142.}\\[1ex]
{\large\bfseries NUMERICAL TABLES SUPPLEMENTARY TO SECOND MEMOIR ON QUANTICS.}
\end{center}

\noindent\textit{[Now first published (1889).]}

\medskip

\noindent In the present paper I arrange in a more compendious form and continue to a much greater extent the tables (first of each pair) given Nos.~35--39 of my Second Memoir on Quantics, 141, pp.~260--264, which relate to the cubic, the quartic and the quintic functions; and I give the like tables for the sextic, the septimic and the octavic functions respectively. The cubic table exhibits the coefficients of the several terms of the function $1 \div (1 - z \cdot 1 - xz \cdot 1 - x^2 z \cdot 1 - x^3 z)$, or, what is the same thing, it gives the number of partitions of a given number into a given number of parts, the parts being $0, 1, 2, 3$ (repetitions admissible): or again, regarding the letters $a, b, c, d$, as having the weights $0, 1, 2, 3$ respectively, it shows the number of literal terms of a given degree and given weight. And similarly for the quartic, quintic, sextic, septimic and octavic tables respectively, the parts of course being $0, 1, \ldots$ up to $4, 5, 6, 7$ or $8$, and the letters being $a, b, \ldots$ up to $e, f, g, h$ or $i$. The extent of the tables is as follows:
\begin{center}
\begin{tabular}{lll}
cubic table & extends to deg-weight & 18--27 \\
quartic     & ,, & 18--36 \\
quintic     & ,, & 18--45 \\
sextic      & ,, & 16--45 \\
septimic    & ,, & 12--42 \\
octavic     & ,, & 10--40
\end{tabular}
\end{center}
viz.\ for the quintic, the sextic and the octavic functions these are the deg-weights of the highest invariants respectively. I designate the Tables as the $ad$-, $ae$-, $af$-, $ag$-, $ah$- and $ai$-tables respectively.

It is to be noticed that in the several tables the lower part of each column is for shortness omitted; the column has to be completed by taking into it the series of bottom terms of each of the preceding columns: thus in the $af$- or quintic table the complete column for degree 3 would be
\begin{center}
\begin{tabular}{cc}
W & D 3 \\
\hline
$-0$ & 6 \\
$-1$ & 6 \\
$-2$ & 5 \\
$-3$ & 4 \\
$-4$ & 3 \\
$-5$ & 2 \\
$-6$ & 1 \\
$-7$ & 1 \\
\hline
$-$ & 2 \\
$-$ & 1 \\
$-$ & 1
\end{tabular}
\end{center}
where the concluding terms $2, 1, 1$ are the bottom terms of the three preceding columns respectively. And the meaning is that for degree 3, and weight $8$, or $7$, the number of terms is $= 6$; for weight $7-1, = 6$, the number of terms is $= 6$; and similarly for weights $5, 4, 3, 2, 1$, the numbers are $5, 4, 3, 2, 1, 1$; the numbers are those of the terms

\textit{[Figure: Eight-column array showing the triangle of literal terms (for the quintic) in $a, b, c, d, e, f$ broken into columns by weight. Each column lists the products $a^3$, $a^2 b$, $a^2 c$, \ldots\ down to terms in $f$, with the count of terms equal to $6, 6, 5, 4, 3, 2, 1, 1$ for weights $8, 7, \ldots, 1$ respectively.]}

The like remarks and explanations apply to the other tables.

\bigskip
\begin{center}\textbf{$ad$-TABLE.}\end{center}

\textit{[Figure: Triangular numerical table giving the coefficients $a_{D, W}$ for the cubic, with columns labelled by degree $D = 0, 1, 2, \ldots, 18$ and rows labelled by weight $W$ from $0$ down to $-9$; the displayed entries reach weight $-27$. The first row (column $a = 1$) reads $1, 3, 6, 4, 6, 8, 7, 9, 11, 10, 13, 14, 13, 16, 17, 16, 18, 20, 19, 21, 23, 22, 24, 26, 25, 27$.]}

\bigskip
\begin{center}\textbf{$ae$-TABLE.}\end{center}

\textit{[Figure: Triangular numerical table for the quartic, with columns labelled by degree $D = 0, 1, 2, \ldots, 18$ and rows labelled by weight $W$ from $0$ down to $-18$; entries reach weight $-36$. Each cell shows the count of literal terms of the indicated degree and weight.]}

\bigskip
\begin{center}\textbf{$af$-TABLE.}\end{center}

\textit{[Figure: Triangular numerical table for the quintic, columns labelled by degree $D = 0, 1, 2, \ldots, 18$ and rows labelled by weight $W$ from $0$ down to $-27$; entries reach weight $-45$.]}

\bigskip
\begin{center}\textbf{$ag$-TABLE.}\end{center}

\textit{[Figure: Triangular numerical table for the sextic, columns labelled by degree $D = 0, 1, 2, \ldots, 15$ and rows labelled by weight $W$ from $0$ down to $-30$; entries reach weight $-45$.]}

\bigskip
\begin{center}\textbf{$ah$-TABLE.}\end{center}

\textit{[Figure: Triangular numerical table for the septimic, columns labelled by degree $D = 0, 1, 2, \ldots, 12$ and rows labelled by weight $W$ from $0$ down to $-30$; entries reach weight $-42$.]}

\bigskip
\begin{center}\textbf{$ai$-TABLE.}\end{center}

\textit{[Figure: Triangular numerical table for the octavic, columns labelled by degree $D = 0, 1, 2, \ldots, 10$ and rows labelled by weight $W$ from $0$ down to $-30$; entries reach weight $-40$.]}

The numbers of each table are connected in several ways with those of the preceding tables. One of these connexions, which is of some importance, is best explained by an example: in the $af$-table, $8$-$20$, the number of terms of degree 8 and weight $20$ is $73$; and we have $73 = 1 + 6 + 16 + 23 + 27$, viz.\ (see p.~278) those are the numbers of the terms in $a^5, a^4, a^3, a^2, a$ respectively: the complementary factors, (for example) of $a^4$ are $bef^2$, \&c.\ terms in $b, c, d, e, f$, of the degree 5 and weight $20$, and (replacing therein each letter by that which immediately precedes it) these are in number equal to the terms in $a, b, c, d, e$ of the degree 5 and weight $20 - 5, = 15$; thus the number 6 of the terms in question is that for the deg-weight $5$-$15$ of the $ae$-table: and so $1, 6, 16, 23, 27$ are the numbers in the $ae$-table for the deg-weights $4$-$16, 5$-$15, 6$-$14, 7$-$13$ and $8$-$12$ respectively, or (making a change rendered necessary by the abbreviated form of the tables) say for the deg-weights $4$-$0, 5$-$15, 6$-$14, 7$-$13$ and $8$-$12$.

%==============================================================
% Paper 143
%==============================================================

\bigskip\bigskip
\begin{center}
{\Large\bfseries 143.}\\[1ex]
{\large\bfseries TABLES OF THE COVARIANTS M TO W OF THE BINARY QUINTIC: FROM THE SECOND, THIRD, FIFTH, EIGHTH, NINTH AND TENTH MEMOIRS ON QUANTICS.}
\end{center}

\noindent\textit{[Arranged in the present form, 1889.]}

\medskip

\noindent The binary quintic has in all (including the quintic itself and the invariants) 23 covariants, which I have represented by the capital letters, $A, B, C, \ldots, W$ (alternative forms of two of these are denoted by $Q'$ and $S'$). The covariants $A, \ldots, L$, and also $Q, Q'$ were given in my Second Memoir on Quantics, and except $Q$ and $Q'$ are reproduced in the present reprint thereof, 141; in all these I gave not only the literal terms actually presenting themselves, but also the terms with zero coefficients; in the other covariants however, or in most of them, the terms with zero coefficients were omitted. It is very desirable to have in every case the complete series of literal terms, and in the covariants as here printed they are accordingly inserted: the number of terms is in each case known beforehand by the foregoing $af$-table, 142, and any omission is thus precluded; by means of this $af$-table we have the numbers of terms as shown in the following list.

I have throughout (as was done in the Ninth and Tenth Memoirs) expressed the literal terms in a slightly different form from that employed in the Second Memoir: this is done in order to show at a glance in each column the set of terms which contain a given power of $a$, and in each such set the terms which contain a given power of $b$.

The numerical verifications are also given not only for the entire column but for each set of terms containing the same power of $a$; viz.\ in most cases, but not always, the positive and negative coefficients of a set have equal sums, which are shown by a number with the sign $\pm$ prefixed. The verification is in some cases given in regard to the subsets involving the same powers of $a$ and $b$, here also the sums of the positive and negative coefficients are not in every case equal. The cases of inequality will be referred to at the end of this paper.

The whole series of covariants is as follows:

\begin{center}
\begin{tabular}{cllll}
Mem. & No.~of table & & & deg-weight \\
\hline
2 & 13 & $A =$ & $(1, 1, 1, 1, 1, 1 \mid x, y)^5$ & 1 (0\ldots 5) \\
,, & 14 & $B =$ & $(3, 3, 3 \mid x, y)^2$ & 2 (4..6) \\
,, & 15 & $C =$ & $(2, 2, 3, 3, 3, 2, 2 \mid x, y)^6$ & 2 (2\ldots 8) \\
,, & 16 & $D =$ & $(6, 6, 6, 8 \mid x, y)^3$ & 3 (6..9) \\
,, & 17 & $E =$ & $(5, 6, 6, 6, 6, 5 \mid x, y)^5$ & 3 (5\ldots 10) \\
,, & 18 & $F =$ & $(3, 4, 5, 6, 6, 6, 6, 5, 4, 3 \mid x, y)^9$ & 3 (3\ldots\ldots 12) \\
,, & 19 & $G =$ & $(12 \mid x, y)^0$, Invt. & 4--10 \\
,, & 20 & $H =$ & $(11, 11, 12, 11, 11 \mid x, y)^4$ & 4 (8\ldots 12) \\
,, & 21 & $I =$ & $(9, 11, 11, 12, 11, 11, 9 \mid x, y)^6$ & 4 (7\ldots\ldots 13) \\
,, & 22 & $J =$ & $(20, 20 \mid x, y)^1$ & 5 (12, 13) \\
,, & 23 & $K =$ & $(19, 20, 20, 19 \mid x, y)^3$ & 5 (11..14) \\
,, & 24 & $L =$ & $(16, 18, 19, 20, 20, 19, 18, 16 \mid x, y)^7$ & 5 (9\ldots\ldots 16) \\
8 & 83 & $M =$ & $(32, 32, 32 \mid x, y)^2$ & 6 (14..16) \\
,, & 84 & $N =$ & $(30, 32, 32, 32, 30 \mid x, y)^4$ & 6 (13\ldots\ldots 17) \\
9 & 90 & $O =$ & $(49, 49 \mid x, y)^1$ & 7 (17, 18) \\
,, & 91 & $P =$ & $(46, 48, 49, 49, 48, 46 \mid x, y)^5$ & 7 (15\ldots\ldots 20) \\
2 & Q 25 & $Q, Q' =$ & $(73 \mid x, y)^0$, Invt. & 8--20 \\
  & Q' 26 & & & \\
9 & 92 & $R =$ & $(71, 73, 71 \mid x, y)^2$ & 8 (19..21) \\
9 & S 93 & $S, S' =$ & $(101, 102, 102, 101 \mid x, y)^3$ & 9 (21\ldots 24) \\
10 & S 93 bis & & & \\
9 & 94 & $T =$ & $(190, 190 \mid x, y)^1$ & 11 (27, 28) \\
3 & 29 & $U =$ & $(252 \mid x, y)^0$, Invt. & 12--30 \\
9 & 95 & $V =$ & $(325, 325 \mid x, y)^1$ & 13 (32, 33) \\
5 & 29a & $W =$ & $(967 \mid x, y)^0$, Invt. & 18--45
\end{tabular}
\end{center}

\textit{[The detailed Table M, No.~83 (and the remaining tables N through W, Nos.~84 through 95 with the bis-tables) lists the literal terms of each covariant in columns by power of $a$, then $b$, with their numerical coefficients. The M-table has a three-column layout showing terms of $a^3 b^3 f^2$, $a^3 b^2 cf^2$, etc., with column-sums $\pm 32$, $\pm 167$, $\pm 52$ for verification of the first set, second set, and third set respectively.]}

\bigskip
\noindent\textit{[The remaining detailed tables of covariants $N$ through $W$ are not reproduced here in full; they are extensive numerical tables that document the literal terms and coefficients of each covariant, with column-sum verifications.]}

\end{document}
